%%%%%%%%%%%%%%%%%%%%%%%%%%%%%%%%%%%%%%%%%%%%%%%%%%%%%%%%%%%%%%%%%%%%%%%%%%%%%%
% Manuscript v3 (P3). TARGET_YEAR = 2027, TEMPLATE_YEAR = 2026, SUBMISSION_READY = false.
% No official ICML 2027 style was found on 2026-09-26 (the URLs listed in paper/README.md
% returned 404); the unmodified official ICML 2026 style files are used temporarily. Do not
% rename the style or change margins/fonts.
% Order (v3): (a) raw vs canonical coordinates -> (b) admissible feature realization ->
% (c) matched-denominator robustness -> (d) lexical change / quality tradeoff.
% Every number in the text comes from generated_numbers.tex or tables/*.tex, which
% src/saeedit/paper_assets.py writes from results/. Missing results are marked
% \todo{EXPERIMENT_ID: required evidence}.
%%%%%%%%%%%%%%%%%%%%%%%%%%%%%%%%%%%%%%%%%%%%%%%%%%%%%%%%%%%%%%%%%%%%%%%%%%%%%%

\documentclass{article}

\usepackage{microtype}
\usepackage{graphicx}
\usepackage{subcaption}
\usepackage{booktabs}
\usepackage{hyperref}

\newcommand{\theHalgorithm}{\arabic{algorithm}}

% Initial blind version for review:
\usepackage{icml2026}

\usepackage{amsmath}
\usepackage{amssymb}
\usepackage{mathtools}
\usepackage{amsthm}
\usepackage[capitalize,noabbrev]{cleveref}
\usepackage{pgfplots}
\usepgfplotslibrary{groupplots}
\pgfplotsset{compat=1.16}

\theoremstyle{plain}
\newtheorem{theorem}{Theorem}[section]
\newtheorem{proposition}[theorem]{Proposition}
\newtheorem{lemma}[theorem]{Lemma}
\newtheorem{corollary}[theorem]{Corollary}
\theoremstyle{definition}
\newtheorem{definition}[theorem]{Definition}
\newtheorem{assumption}[theorem]{Assumption}
\theoremstyle{remark}
\newtheorem{remark}[theorem]{Remark}

% Inline marker for results that do not exist yet (not todonotes: margins are not changed).
\newcommand{\todo}[1]{\textcolor{red}{[TODO\,--\,#1]}}
\newcommand{\notrun}{\textsc{not run}}
\newcommand{\R}{\mathbb{R}}
\newcommand{\LN}{\textsc{ln}}

% Generated by src/saeedit/paper_assets.py from results/raw/*.csv. Do not edit by hand.
\newcommand{\ClAdmMax}{$3.6\times10^{-12}$}
\newcommand{\ClCPU}{1481}
\newcommand{\ClCanonMax}{$3.8\times10^{-15}$}
\newcommand{\ClGroups}{16}
\newcommand{\ClKLMax}{$1.1\times10^{-15}$}
\newcommand{\ClLogitMax}{$2.8\times10^{-13}$}
\newcommand{\ClMeanCanonMax}{$1.4\times10^{-13}$}
\newcommand{\ClMeanKL}{$1.4\times10^{-15}$}
\newcommand{\ClMeanRawMed}{99}
\newcommand{\ClPdeltaNorm}{$3.1\times10^{-15}$}
\newcommand{\ClRawMed}{124.7}
\newcommand{\ClRawRepro}{$4.6\times10^{-16}$}
\newcommand{\ClRows}{760}
\newcommand{\ClStoredRepro}{0}
\newcommand{\ClUnprojMax}{0.34}
\newcommand{\ClWall}{843}
\newcommand{\CxActRel}{$1.0\times10^{-14}$}
\newcommand{\CxActRelUnc}{0.0319}
\newcommand{\CxEightMax}{1.14}
\newcommand{\CxEightMeanMax}{6}
\newcommand{\CxEightMeanMin}{3}
\newcommand{\CxEightMed}{0.79}
\newcommand{\CxEightMin}{0.51}
\newcommand{\CxFSixFourKL}{$1.5\times10^{-15}$}
\newcommand{\CxFSixFourProb}{$5.1\times10^{-14}$}
\newcommand{\CxFThreeTwoKL}{$7.5\times10^{-7}$}
\newcommand{\CxFThreeTwoProb}{$2.8\times10^{-5}$}
\newcommand{\CxFVE}{0.83}
\newcommand{\CxFVEUnc}{0.003}
\newcommand{\CxLZero}{72.4}
\newcommand{\CxLZeroUnc}{176}
\newcommand{\CxMeanDrift}{125}
\newcommand{\CxMeanKLSixFour}{$1.3\times10^{-15}$}
\newcommand{\CxMeanKLThreeTwo}{$6.0\times10^{-7}$}
\newcommand{\CxMeanMaxFeat}{37.5}
\newcommand{\CxMedNorm}{101}
\newcommand{\CxRho}{10.1}
\newcommand{\CxShift}{279}
\newcommand{\DenseJumpActMax}{23}
\newcommand{\DenseJumpActMed}{11}
\newcommand{\DenseJumpActMin}{3}
\newcommand{\DenseReluActFourWins}{0/20}
\newcommand{\DenseReluActMax}{22}
\newcommand{\DenseReluActMed}{13}
\newcommand{\DenseReluActMin}{5}
\newcommand{\DenseReluActOneWins}{3/20}
\newcommand{\DenseReluActQWins}{10/20}
\newcommand{\DenseReluInactFourWins}{0/20}
\newcommand{\DenseReluInactOneWins}{0/20}
\newcommand{\DenseReluInactQWins}{0/20}
\newcommand{\EncGradLeakRzero}{0.283}
\newcommand{\FeasTotal}{260}
\newcommand{\FeasTotalPass}{260}
\newcommand{\LxCalSecVOne}{158}
\newcommand{\LxCalSecVOneOne}{630}
\newcommand{\LxHitsVOne}{5}
\newcommand{\LxHitsVOneOne}{13}
\newcommand{\LxStatus}{BLOCKED}
\newcommand{\LxWindowsAll}{1943}
\newcommand{\PilotAllDec}{0.656}
\newcommand{\PilotAllDecInact}{0.716}
\newcommand{\PilotAllEnc}{0.778}
\newcommand{\PilotAllEncInact}{1.36}
\newcommand{\PilotAllLn}{0.114}
\newcommand{\PilotAllLnInact}{0.295}
\newcommand{\PilotAllMean}{52.5}
\newcommand{\PilotAllMeanInact}{184}
\newcommand{\PilotAllRand}{108}
\newcommand{\PilotAllRandInact}{348}
\newcommand{\PilotDAllAct}{-0.503\,\allowbreak[-0.584, -0.421]}
\newcommand{\PilotDAllInact}{-0.43\,\allowbreak[-0.517, -0.35]}
\newcommand{\PilotDKLAct}{-0.00377\,\allowbreak[-0.0117, -0.00025]}
\newcommand{\PilotDKLInact}{-0.00234\,\allowbreak[-0.00651, -0.00087]}
\newcommand{\PilotDNewAct}{-0.17\,\allowbreak[-0.202, -0.142]}
\newcommand{\PilotDNewInact}{-0.0557\,\allowbreak[-0.098, -0.00262]}
\newcommand{\PilotDocs}{54}
\newcommand{\PilotEqAllDec}{0.517}
\newcommand{\PilotEqAllEnc}{1.16}
\newcommand{\PilotEqAllLn}{0.142}
\newcommand{\PilotEqAllMean}{116}
\newcommand{\PilotEqAllRand}{0.341}
\newcommand{\PilotFeatures}{8}
\newcommand{\PilotFwdSec}{22.6}
\newcommand{\PilotGOne}{continue side}
\newcommand{\PilotGThree}{continue side}
\newcommand{\PilotGTwo}{continue side}
\newcommand{\PilotGainDec}{0.895}
\newcommand{\PilotGainDecInact}{$\approx 0$}
\newcommand{\PilotGainEnc}{1.5}
\newcommand{\PilotGainEncInact}{0.433}
\newcommand{\PilotGainLn}{1.35}
\newcommand{\PilotGainLnInact}{0.304}
\newcommand{\PilotGainMean}{1.85}
\newcommand{\PilotGainMeanInact}{0.581}
\newcommand{\PilotGainRand}{0.0383}
\newcommand{\PilotGainRandInact}{$\approx 0$}
\newcommand{\PilotKLDec}{5.32}
\newcommand{\PilotKLEnc}{1.03}
\newcommand{\PilotKLLn}{0.96}
\newcommand{\PilotKLMean}{$\approx 0$}
\newcommand{\PilotKLRand}{691}
\newcommand{\PilotLnEncAll}{-0.687\,\allowbreak[-0.822, -0.613]}
\newcommand{\PilotLnEncKL}{-0.000058\,\allowbreak[-0.00017, -0.000012]}
\newcommand{\PilotLnEncNorm}{0.99\,\allowbreak[0.895, 1.12]}
\newcommand{\PilotMeanInfAct}{0/64}
\newcommand{\PilotMeanInfInact}{0/64}
\newcommand{\PilotMeanKLAct}{\ensuremath{-7.0\times10^{-18}}\,\allowbreak[\ensuremath{-4.5\times10^{-17}}, \ensuremath{4.2\times10^{-17}}]}
\newcommand{\PilotMeanKLInact}{\ensuremath{2.3\times10^{-17}}\,\allowbreak[\ensuremath{-3.4\times10^{-17}}, \ensuremath{7.1\times10^{-17}}]}
\newcommand{\PilotNewDec}{0.311}
\newcommand{\PilotNewEnc}{0.366}
\newcommand{\PilotNewLn}{0.114}
\newcommand{\PilotNewMean}{52.1}
\newcommand{\PilotNewRand}{108}
\newcommand{\PilotNormDec}{21.6}
\newcommand{\PilotNormEnc}{11.8}
\newcommand{\PilotNormLn}{12.7}
\newcommand{\PilotNormMean}{9.04}
\newcommand{\PilotNormRand}{438}
\newcommand{\PilotOrigDec}{0.583}
\newcommand{\PilotOrigLn}{$\approx 0$}
\newcommand{\PilotQualityGroups}{32}
\newcommand{\PilotRandInf}{272}
\newcommand{\PilotRandTotal}{512}
\newcommand{\PilotRescInactTgt}{0.727}
\newcommand{\PilotRows}{6144}
\newcommand{\PilotSolveMs}{22}
\newcommand{\RepairConvRJMax}{80}
\newcommand{\RepairConvRJMin}{72}
\newcommand{\RepairConvTKMax}{66}
\newcommand{\RepairConvTKMin}{2}
\newcommand{\RepairRowsPerCell}{80}
\newcommand{\RepairTgtMax}{0.84}
\newcommand{\RepairTgtMin}{0.23}
\newcommand{\RfiveDecTgt}{0.669}
\newcommand{\RfiveLnLeak}{2.21}
\newcommand{\RfiveLnNorm}{3.6}
\newcommand{\RfiveRescLeak}{1.44}
\newcommand{\RfiveRescNorm}{3.1}
\newcommand{\RfouraLnTgt}{0.675}
\newcommand{\RfouraRes}{0.821}
\newcommand{\RfouraResMin}{0.784}
\newcommand{\RfouraResRowMin}{0.696}
\newcommand{\RfourbDecLeak}{1.73}
\newcommand{\RfourbDecLeakP}{0.638}
\newcommand{\RfourbDecLeakU}{1.59}
\newcommand{\RfourbLnLeak}{2.16}
\newcommand{\RfourbLnNorm}{1.35}
\newcommand{\RgAllBelowHundred}{35}
\newcommand{\RgAllNPosMin}{8}
\newcommand{\RgGOneAct}{0.583\,\allowbreak[0.452, 0.6]}
\newcommand{\RgGOneActN}{42}
\newcommand{\RgGOneAsRun}{0.426\,\allowbreak[0.375, 0.469]}
\newcommand{\RgGOneInact}{0.342\,\allowbreak[0.289, 0.371]}
\newcommand{\RgGOneInactN}{40}
\newcommand{\RgInf}{272}
\newcommand{\RgNPosMax}{134}
\newcommand{\RgNPosMin}{26}
\newcommand{\RgOK}{4848}
\newcommand{\RgPrimExclMax}{0\%}
\newcommand{\RgRawOnly}{512}
\newcommand{\RgRescActWithin}{256}
\newcommand{\RgRescInactErr}{0.727}
\newcommand{\RgRescInactWithin}{0}
\newcommand{\RgSecExclMax}{55\%}
\newcommand{\RgSecExclMin}{52\%}
\newcommand{\RgTgtErrMax}{$5.1\times10^{-8}$}
\newcommand{\RgTopTwo}{3}
\newcommand{\RgUnm}{1024}
\newcommand{\RoneDecTgt}{0.365}
\newcommand{\RthreeLnNorm}{1.24}
\newcommand{\RthreeRescLeak}{0.989}
\newcommand{\RtwoDecLeak}{0.794}
\newcommand{\RtwoLnNorm}{1.14}
\newcommand{\RtwoTrLeak}{0.0305}
\newcommand{\RtwoTrTgt}{0.111}
\newcommand{\SpAttempted}{32}
\newcommand{\SpBoth}{1}
\newcommand{\SpBound}{10.9}
\newcommand{\SpBoundOne}{8.9}
\newcommand{\SpBoundTwo}{10.9}
\newcommand{\SpBudDec}{0.1}
\newcommand{\SpBudDm}{0.1}
\newcommand{\SpBudEnc}{0.1}
\newcommand{\SpBudLn}{0.1}
\newcommand{\SpCalDocs}{60}
\newcommand{\SpCalSec}{961}
\newcommand{\SpCalStatus}{RAN}
\newcommand{\SpChangedDec}{26}
\newcommand{\SpChangedDm}{31}
\newcommand{\SpChangedEnc}{24}
\newcommand{\SpChangedLn}{26}
\newcommand{\SpChangedMean}{0}
\newcommand{\SpCtrlKL}{$1.2\times10^{-16}$}
\newcommand{\SpCtrlSame}{32/32}
\newcommand{\SpDajDec}{1.73}
\newcommand{\SpDajDm}{0.00}
\newcommand{\SpDajEnc}{3.92}
\newcommand{\SpDajLn}{2.99}
\newcommand{\SpDajMean}{0.00}
\newcommand{\SpDajNone}{0.00}
\newcommand{\SpDallDec}{3.65}
\newcommand{\SpDallDm}{3.56}
\newcommand{\SpDallEnc}{9.17}
\newcommand{\SpDallLn}{1.68}
\newcommand{\SpDallMean}{0.00}
\newcommand{\SpDallNone}{0.00}
\newcommand{\SpDiscA}{0}
\newcommand{\SpDiscB}{0}
\newcommand{\SpDiscPrim}{0}
\newcommand{\SpDmCPHi}{25.0}
\newcommand{\SpDmCPLo}{2.0}
\newcommand{\SpDmDisc}{3}
\newcommand{\SpDmGain}{2}
\newcommand{\SpDmLoss}{1}
\newcommand{\SpDmNet}{+1}
\newcommand{\SpDocs}{32}
\newcommand{\SpFailed}{0}
\newcommand{\SpFeature}{4671}
\newcommand{\SpHitDec}{3 [0, 9]}
\newcommand{\SpHitDm}{6 [0, 16]}
\newcommand{\SpHitEnc}{3 [0, 9]}
\newcommand{\SpHitLn}{3 [0, 9]}
\newcommand{\SpHitMean}{3 [0, 9]}
\newcommand{\SpHitNone}{3 [0, 9]}
\newcommand{\SpKLDec}{0.0281}
\newcommand{\SpKLDm}{0.0487}
\newcommand{\SpKLEnc}{0.0269}
\newcommand{\SpKLLn}{0.0205}
\newcommand{\SpKLMean}{$1.5\times10^{-17}$}
\newcommand{\SpKLNone}{0}
\newcommand{\SpKappa}{0.1}
\newcommand{\SpMaxDoc}{9}
\newcommand{\SpMedNorm}{101.6}
\newcommand{\SpNLLDec}{1.20}
\newcommand{\SpNLLDm}{1.34}
\newcommand{\SpNLLEnc}{1.27}
\newcommand{\SpNLLLn}{1.37}
\newcommand{\SpNLLMean}{1.32}
\newcommand{\SpNLLNone}{1.32}
\newcommand{\SpNeither}{31}
\newcommand{\SpOverlapCal}{0}
\newcommand{\SpOverlapPilot}{28}
\newcommand{\SpPosDocs}{22}
\newcommand{\SpPosWin}{58}
\newcommand{\SpPrim}{0.0}
\newcommand{\SpPrimHi}{0.0}
\newcommand{\SpPrimLo}{0.0}
\newcommand{\SpRunSec}{1334}
\newcommand{\SpRunStatus}{RAN}
\newcommand{\SpSkipped}{0}
\newcommand{\SpStarWin}{28}
\newcommand{\SpVerdict}{not supported}
\newcommand{\SpWin}{1943}
\newcommand{\SparseFiftyActFourWins}{4/20}
\newcommand{\SparseFiftyActMax}{11}
\newcommand{\SparseFiftyActMed}{4}
\newcommand{\SparseFiftyActOneWins}{19/20}
\newcommand{\SparseFiftyActQWins}{20/20}
\newcommand{\SparseFiftyInactFourWins}{3/20}
\newcommand{\SparseFiftyInactOneWins}{11/20}
\newcommand{\SparseFiftyInactQWins}{16/20}
\newcommand{\SparseSeventyFiveActFourWins}{17/20}
\newcommand{\SparseSeventyFiveActMax}{6}
\newcommand{\SparseSeventyFiveActMed}{1}
\newcommand{\SparseSeventyFiveActOneWins}{18/20}
\newcommand{\SparseSeventyFiveActQWins}{19/20}
\newcommand{\SparseSeventyFiveInactFourWins}{15/20}
\newcommand{\SparseSeventyFiveInactOneWins}{17/20}
\newcommand{\SparseSeventyFiveInactQWins}{18/20}
\newcommand{\TrWinsMax}{20}
\newcommand{\TrWinsMin}{10}
\newcommand{\TxBothDecExt}{0.223}
\newcommand{\TxBothDecInt}{0.313}
\newcommand{\TxBothDiff}{-0.00137}
\newcommand{\TxBothLnExt}{0.223}
\newcommand{\TxBothLnInt}{$<\!10^{-6}$}
\newcommand{\TxBothWins}{11/20}
\newcommand{\TxDecAccDecExt}{$<\!10^{-6}$}
\newcommand{\TxDecAccDecInt}{0.242}
\newcommand{\TxDecAccDiff}{+0.239}
\newcommand{\TxDecAccLnExt}{0.239}
\newcommand{\TxDecAccLnInt}{$<\!10^{-6}$}
\newcommand{\TxDecAccWins}{0/20}
\newcommand{\TxEncAccDecExt}{0.221}
\newcommand{\TxEncAccDecInt}{0.221}
\newcommand{\TxEncAccDiff}{-0.221}
\newcommand{\TxEncAccLnExt}{$<\!10^{-6}$}
\newcommand{\TxEncAccLnInt}{$<\!10^{-6}$}
\newcommand{\TxEncAccWins}{20/20}
\newcommand{\VfiveMaxDiff}{$7.5\times10^{-13}$}
\newcommand{\VsixMaxErr}{3.3e-15}
\newcommand{\VsixRows}{550}


\icmltitlerunning{Auditing SAE Feature Edits in GPT-2 Small}

\begin{document}

\twocolumn[
  \icmltitle{Auditing Sparse Autoencoder Feature Edits in GPT-2 Small:\\
    Coordinates, Fidelity, and a Lexical Pilot}

  \begin{icmlauthorlist}
    \icmlauthor{Anonymous Author}{anon}
  \end{icmlauthorlist}

  \icmlaffiliation{anon}{Anonymous Institution}
  \icmlcorrespondingauthor{Anonymous Author}{anonymous@example.com}

  \icmlkeywords{sparse autoencoders, steering, interpretability, interventions, evaluation}

  \vskip 0.3in
]

\printAffiliationsAndNotice{}

\begin{abstract}
  Steering with a sparse autoencoder (SAE) edits an activation along a decoder direction and
  reports the edit as a change of one feature. For GPT-2 small and a public layer-8
  residual SAE we audit that report in four steps and propose no method.
  (a) \emph{Coordinates.} We prove, for GPT-2 only, that edits along the all-ones direction
  leave the logits unchanged, and we read edits canonically as $E(P(h+\delta))$. A drift of
  \CxMeanDrift{} times the edit norm that we had attributed to such a model-invisible edit
  comes from the raw readout $E(x+\delta)$; canonically it is at most \ClCanonMax. Our fast
  implementation matches TransformerLens/SAELens in float64, and the pre-registered float32
  check fails through rounding.
  (b) \emph{Realization.} In a conditional pilot (\PilotFeatures{} features, \PilotDocs{}
  held-out documents), target-matched projected-decoder edits leave non-target drift of
  \PilotAllDec$\alpha$, and a least-norm correction posed with $J_E(Ph)P$ leaves
  \PilotAllLn$\alpha$. In dense toy regimes the same correction moves drift onto newly
  activated features.
  (c) \emph{Denominators.} The main methods are matched at every test point, and an earlier
  readout that mixed matched and unmatched rows is corrected.
  (d) \emph{Behavior.} Our only behavioral test, wedding-lexicon occurrence, is a lexical
  proxy and not semantic topic control. It was blocked at calibration, because the
  calibration text had \LxHitsVOneOne{} keyword windows against a pre-set minimum of 20. We
  report no behavioral effect.
\end{abstract}

%%%%%%%%%%%%%%%%%%%%%%%%%%%%%%%%%%%%%%%%%%%%%%%%%%%%%%%%%%%%%%%%%%%%%%%%%%%%%%
\section{Introduction}
\label{sec:intro}

Sparse autoencoders (SAEs) are used to steer language models by editing an internal
activation along a learned feature. Common implementations add a multiple of the feature's
decoder vector to the residual stream, or clamp the feature inside the SAE reconstruction
while keeping the reconstruction error
\citep{templeton2024scaling,obrien2025steering,wu2025axbench}. The edit is then reported in
the SAE's vocabulary: ``feature $j$ was increased by $\alpha$''. We ask what that sentence
certifies for one public model--SAE pair, GPT-2 small with a residual SAE at layer 8, and we
check its preconditions in a fixed order.
\begin{enumerate}
  \item[(a)] \emph{Coordinates.} Which activation does the SAE read after an edit? We
  separate a \emph{raw off-slice} readout $E(x+\delta)-E(x)$ from the \emph{canonical
  reprojected} readout $E(P(h+\delta))-E(Ph)$, where $x = Ph$ is the SAE's training
  coordinate (\cref{sec:coord}).
  \item[(b)] \emph{Admissible feature realization.} Among edits the model can see, how much
  of the re-encoded change lands on the target, and does a local correction move the rest
  elsewhere (\cref{sec:fidelity})?
  \item[(c)] \emph{Matched denominators.} Are methods compared on the same test points once
  infeasible and unmatched rows are counted (\cref{sec:denominators})?
  \item[(d)] \emph{Lexical change and quality.} Does a better internal realization give more
  wedding-lexicon occurrence at matched output change (\cref{sec:external})? This is a
  lexical proxy, not semantic topic control.
\end{enumerate}

Prior work covers pieces of this. Encoder and decoder directions differ
\citep{chanin2025absorption,oneill2025compute}, and encoder--decoder misalignment has been
named as one cause of steering failure, next to off-manifold artifacts
\citep{luo2026learning}. SAE-targeted steering solves for vectors that hit SAE features
through a fitted effect map, and reports that a pseudo-inverse of that map did worse
\citep{chalnev2024improving}. Least-norm readout edits and null-space protection are
standard \citep{geiger2024finding,wu2024reft,fang2025alphaedit,cui2026sae}. We therefore do
not propose a method.

\paragraph{Contributions.}
\begin{itemize}
  \item \emph{Readout convention} (\cref{sec:coord}). For GPT-2 we prove that adding
  $c\mathbf{1}$ at the hook leaves the logits unchanged, and we define the canonical readout,
  which depends only on the model-visible part of an edit. An executed contract shows that
  our fast path equals the canonical TransformerLens/SAELens path in float64; the
  pre-registered float32 check fails through rounding, and we keep that failure. A
  previously reported ``$\CxMeanDrift\times$ SAE drift'' of a model-invisible edit is a raw
  off-slice readout; the canonical readout of the same edit is at most \ClCanonMax.
  \item \emph{Admissible realization, conditional pilot} (\cref{sec:fidelity}). On
  \PilotDocs{} held-out documents and \PilotFeatures{} features, a least-norm correction posed
  with $J_{\mathrm{eff}} = J_E(Ph)P$ lowers same-layer non-target drift relative to the
  projected decoder at matched target change. Controlled toy regimes show where the same
  correction instead moves drift onto newly activated features.
  \item \emph{Denominators} (\cref{sec:denominators}). The three main methods are compared
  on all test points; random directions lose about half of them, unmatched baselines and
  raw mean-only rows are reported separately, and one earlier gate readout mixed matched and
  unmatched rows.
  \item \emph{Lexical pilot} (\cref{sec:external}). The only behavioral test, wedding-lexicon
  occurrence, was frozen in advance and stopped at calibration: \LxHitsVOneOne{} of
  \LxWindowsAll{} calibration windows contain a keyword, below the pre-set minimum of 20. No
  behavioral number exists.
\end{itemize}
We make no claim about semantic steering, and lower internal drift or lower KL is never
taken as evidence that a method is better.

%%%%%%%%%%%%%%%%%%%%%%%%%%%%%%%%%%%%%%%%%%%%%%%%%%%%%%%%%%%%%%%%%%%%%%%%%%%%%%
\section{Related Work}
\label{sec:related}

\paragraph{SAE steering and its evaluation.}
Feature steering by clamping or adding decoder directions was demonstrated at scale
\citep{templeton2024scaling} and applied to refusal \citep{obrien2025steering}. AxBench
steers with $h + \alpha\, w_{\mathrm{dec}}$, scores concept, instruction relevance and
fluency with a judge, and reports a mean steering score of 0.165 for SAEs against 0.239 for
difference-in-means \citep{wu2025axbench}; we therefore keep a difference-in-means baseline.
Input features and output features rarely coincide \citep{arad2025saes}. Other work edits in
the SAE latent space and decodes with the error term added back \citep{bayat2025steering},
fits decoder-space edits to counterfactual activations \citep{makelov2024towards}, or scores
encoder--decoder interventions by output effects \citep{bhalla2025towards}. These works
measure behavior; they do not re-encode the edited activation.

\paragraph{Targeting SAE features.}
SAE-TS learns a linear map $\hat y = xM + b$ from steering vectors to the SAE activations of
steered rollouts and steers with $M_j/\|M_j\| - \lambda Mb/\|Mb\|$; its appendix reports that
the pseudo-inverse of $M$ is worse on all nine tasks \citep{chalnev2024improving}.
Feature-guided activation additions reuse this map \citep{soo2025interpretable}. We instead
re-encode each edited activation at the edit layer with the exact per-input encoder rows and
control the edit norm. The pseudo-inverse result is prior evidence against our correction.

\paragraph{Encoders are not inverses of decoders.}
Absorption leaves holes in encoder recall \citep{chanin2025absorption}, a linear--nonlinear
encoder has an amortisation gap \citep{oneill2025compute}, and encoding a bare steering
vector is dominated by the encoder bias \citep{mayne2024can}. \citet{luo2026learning}
conclude that off-manifold artifacts, ``not just encoder-decoder misalignment'', contribute
to steering failures. Side effects have been predicted from decoder geometry and logit
effects \citep{duan2026preintervention,khan2026look}; we measure the target's own same-layer
readout instead.

\paragraph{Least-norm and null-space interventions.}
Distributed interchange interventions \citep{geiger2024finding} and LoReFT
\citep{wu2024reft} are least-norm edits that set orthonormal readouts to target values; our
correction has the same form with non-orthogonal SAE encoder rows. AlphaEdit and AlphaSteer
project updates onto null spaces of preserved keys or benign activations
\citep{fang2025alphaedit,sheng2026alphasteer}; LEACE and MiMiC are least-change edits
\citep{belrose2025leace,singh2025representation}. \citet{cui2026sae} hold SAE features fixed
with within- and across-layer null-space projectors while changing behavior, which is our
protected-set constraint used adversarially. Readout and intervention directions are related
through a metric \citep{park2024linear}, so the Euclidean norm is one choice among several.

\paragraph{Benchmarks.}
SAEBench \citep{karvonen2025saebench} and RAVEL
\citep{huang2024ravel,chaudhary2024evaluating} evaluate interventions downstream; on our
reading, none compares intended with re-encoded change. Jacobian SAEs
\citep{farnik2025jacobian} reduce, with an identity in place of the MLP, to the
$E_{\mathrm{active}}D_{\mathrm{active}}$ matrix of our linear analysis.

%%%%%%%%%%%%%%%%%%%%%%%%%%%%%%%%%%%%%%%%%%%%%%%%%%%%%%%%%%%%%%%%%%%%%%%%%%%%%%
\section{Problem Setup}
\label{sec:setup}

\paragraph{SAE.}
Let $x\in\R^d$ be an activation, written in the coordinates the SAE was trained on. The
encoder has pre-activations $p(x) = Ex + b$ with $E\in\R^{m\times d}$ and features
$a = \sigma(p)$, where $\sigma$ is ReLU, JumpReLU with threshold $\theta$, or TopK. The decoder
$D\in\R^{d\times m}$ has columns $D_j$. For the public SAE in \cref{sec:coord},
$p(x) = W_{\mathrm{enc}}^{\top}(x - b_{\mathrm{dec}}) + b_{\mathrm{enc}}$, so
$E = W_{\mathrm{enc}}^{\top}$ and $b = b_{\mathrm{enc}} - W_{\mathrm{enc}}^{\top}b_{\mathrm{dec}}$.

\paragraph{Edit problem.}
An edit is specified by a target $j$, an intended post-activation change $\alpha>0$, and a
protected set $Q$, which we take to be the non-target features active at $x$. The target
value $a_j+\alpha$ requires the pre-activation change
$\Delta p_j^{\mathrm{req}} = a_j + \alpha - p_j$ whenever the target ends up active. The
achieved change of an edit $\delta$ is $\Delta a = \sigma(p(x+\delta)) - \sigma(p(x))$.

\begin{definition}[Edit-fidelity decomposition]
\label{def:fidelity}
Let $A_0$ and $A_1$ be the active sets before and after the edit. The \emph{target error} is
$|\Delta a_j - \alpha|/\alpha$. The drift terms are $\ell_2$ norms of $\Delta a$ over three
disjoint groups of non-target features $i\neq j$:
\begin{itemize}
  \item $D_{\mathrm{kept}}$ over $i\in A_0\cap A_1$;
  \item $D_{\mathrm{deact}}$ over $i\in A_0\setminus A_1$;
  \item $D_{\mathrm{new}}$ over $i\in A_1\setminus A_0$.
\end{itemize}
Originally active drift is $D_{\mathrm{orig}} = (D_{\mathrm{kept}}^2 + D_{\mathrm{deact}}^2)^{1/2}$,
and $D_{\mathrm{all}}$ is the norm over all $i\neq j$. All terms are reported relative to a
scale: $\alpha$ for target-matched comparisons and the norm budget $\rho$ for equal-norm
comparisons.
\end{definition}
For ReLU, JumpReLU and TopK, a feature inactive before and after has zero change. Hence
$D_{\mathrm{all}}^2 = D_{\mathrm{orig}}^2 + D_{\mathrm{new}}^2$. With $Q = A_0\setminus\{j\}$,
$D_{\mathrm{orig}}$ and $D_{\mathrm{new}}$ are exactly the protected and unprotected leakage of
our toy study (checked numerically in our test suite).

\paragraph{Edits compared.}
\begin{itemize}
  \item \emph{Decoder}: $\delta = cD_j$ with $c = \alpha$.
  \item \emph{Rescaled decoder}: $c$ chosen so that $\Delta p_j = \Delta p_j^{\mathrm{req}}$.
  In the toy study $c$ is computed per input. In the real protocol it is the weight-only
  $\alpha/(E_j D_j)$, which is deployable but not matched.
  \item \emph{Encoder gradient}: $\delta\propto E_j^{\top}$, the least-norm edit for the
  target row alone.
  \item \emph{\LN{} correction}: $\delta = M^{+}t$ with $M = E_{S\cup Q}$ (the rows of $E$ for
  the target set $S=\{j\}$ and the protected set $Q$) and
  $t = (\Delta p_j^{\mathrm{req}},0,\dots,0)$.
  \item \emph{\LN{} + repair}: up to three rounds that add newly crossed features to $Q$ as
  equality constraints.
  \item \emph{Trust region}: $\arg\min_{\|\delta\|\le\tau}\|M\delta - t\|$.
  \item \emph{Random}: a random direction of the same norm, as a control.
\end{itemize}

\paragraph{Admissible edits and readouts.}
For GPT-2, adding $c\mathbf{1}$ to the residual at the hook leaves the logits unchanged
(\cref{prop:admissible}), so a model-visible edit is the class $\delta+\operatorname{span}
(\mathbf{1})$. We represent it by $P\delta$ with $P = I - \mathbf{1}\mathbf{1}^{\top}/d$,
the member in $\operatorname{range}(P)$, where the SAE's inputs $x = Ph$ lie. This is a
choice of representative, not the only valid space. We read an edit \emph{canonically},
$\Delta a = \sigma(p(P(h+\delta))) - \sigma(p(Ph))$, which depends only on the class. Every
actor edit below lies in $\operatorname{range}(P)$, so for it the raw and canonical readouts
coincide. The \LN{} correction linearizes the canonical readout, whose Jacobian is
$J_{\mathrm{eff}} = J_E(Ph)\,P$; it therefore solves $(MP)\delta = t$. A \emph{mean-only}
edit along $\mathbf{1}$ is a control, never a method.

\paragraph{Comparison modes and units.}
In an \emph{equal-norm} comparison, every direction is scaled to the same $\|\delta\|$. In a
\emph{target-matched} comparison, every direction is scaled so that $\Delta a_j = \alpha$.
Rows where no non-negative scale achieves this are kept as infeasible. The independent unit
is a random toy instance in \cref{sec:toy} and a held-out document in \cref{sec:pilot}.
Features, positions and base points within a unit are averaged before any summary.

%%%%%%%%%%%%%%%%%%%%%%%%%%%%%%%%%%%%%%%%%%%%%%%%%%%%%%%%%%%%%%%%%%%%%%%%%%%%%%
\section{Raw and Canonical Coordinates}
\label{sec:coord}

A fidelity number means something only if the SAE reads the activations it was trained on
and the edit being scored is the edit the model receives. We executed a contract for one
public pair before measuring anything.
\begin{itemize}
  \item \emph{Model and SAE:} GPT-2 small and the residual SAE at
  \texttt{blocks.8.hook\_resid\_pre} ($d=768$, $m=24576$), both at pinned revisions.
  \item \emph{Canonical path:} TransformerLens with every processing option passed
  explicitly, and the SAE loaded through SAELens. \emph{Fast path:} the plain Hugging Face
  model used by our adapter. Both receive the same token ids (BOS prepended, 128-token
  windows).
\end{itemize}

\begin{proposition}[Mean invariance, GPT-2 only]
\label{prop:admissible}
Consider the GPT-2 architecture in evaluation mode and exact arithmetic: each block reads
the residual stream only through \texttt{ln\_1} and \texttt{ln\_2}, the unembedding reads
it only through \texttt{ln\_f}, and each LayerNorm subtracts the mean over the $d$ model
coordinates and divides by the standard deviation of the centred vector. Adding
$c_t\mathbf{1}$ to the input of block $\ell$ at any set of positions $t$ leaves every logit
unchanged. We do not claim this for other architectures; for example, RMSNorm does not
subtract the mean, and there the statement is unproved and generally false.
\end{proposition}
The proof is in \cref{app:proofs}. Two consequences shape the audit.
\begin{itemize}
  \item \emph{Canonical readout.} The SAE was trained with
  \texttt{center\_writing\_weights}, so its input is $x = Ph$, where $h$ is the Hugging Face
  hidden state. The canonical readout $E(P(h+\delta)) - E(Ph)$ depends only on the class
  $\delta + \operatorname{span}(\mathbf{1})$. The \emph{raw off-slice} readout
  $E(x+\delta) - E(x)$ feeds the SAE an input outside $\operatorname{range}(P)$ whenever
  $\mathbf{1}^{\top}\delta\neq0$, and then depends on a component the model ignores. The two
  coincide on $\operatorname{range}(P)$.
  \item \emph{Effective Jacobian.} Linearizing the canonical readout gives
  $J_{\mathrm{eff}} = J_E(Ph)\,P$, with $J_{\mathrm{eff}}\mathbf{1} = 0$. The least-norm
  solution of $(MP)\delta = t$ lies in
  $\operatorname{range}(PM^{\top})\subseteq\operatorname{range}(P)$ and spends no norm on
  $\mathbf{1}$; the solution of $M\delta = t$ generally does.
\end{itemize}

\paragraph{Contract execution.}
\Cref{tab:contract-exec} lists the checks and tolerances fixed before execution.
\begin{itemize}
  \item \emph{float32 (pre-registered) fails.} Next-token probabilities of the two paths
  differ by up to \CxFThreeTwoProb, above the $10^{-5}$ tolerance, and the mean-only edit
  moves the distribution by a KL of up to \CxMeanKLThreeTwo, above its $10^{-8}$ bound. We
  keep this failure.
  \item \emph{float64 diagnostic passes.} A rerun designed after the float32 failure, with
  the same tolerances, passes every check with differences of order $10^{-13}$ to
  $10^{-15}$. The float32 discrepancies are rounding, amplified by a common logit shift of
  up to \CxShift{} per position from \texttt{center\_unembed}; all real-model measurements
  use float64.
\end{itemize}

\paragraph{Findings about coordinates.}
\begin{itemize}
  \item \emph{Centring matters.} The uncentred hidden state gives L0~\CxLZeroUnc{} and
  FVE~\CxFVEUnc, against \CxLZero{} and \CxFVE{} in the training coordinates.
  \item \emph{Unprojected edits make the readout convention-dependent.} A raw decoder row
  has a mean component of only \CxEightMeanMin--\CxEightMeanMax\% of its norm, yet its raw
  and canonical readouts differ by \CxEightMin--\CxEightMax{} times the edit's own readout
  change, while the model's log-probabilities for the two representatives agree to
  \ClLogitMax.
  \item \emph{The mean-only ``drift'' is a raw off-slice readout (readout closure).} A
  mean-only edit of norm \CxRho{} (10\% of the median residual norm) leaves the model
  unchanged (KL at most \CxMeanKLSixFour). Read raw, as in the contract's K9 check, it
  moves the SAE by \CxMeanDrift{} times its norm, which we first read as internal change
  without external effect.
  Recomputing the stored contract cases shows that this number is the raw readout
  $E(x_0+\delta)$: we reproduce it to a relative \ClRawRepro, while the canonical readout
  of the same edits is at most \ClCanonMax{} times the norm. On \ClGroups{} pilot test points
  (\ClRows{} edits), raw and canonical readouts of every admissible edit agree to \ClAdmMax,
  and the stored pilot values are reproduced exactly. The number is therefore an
  off-training-subspace stress test of the encoder, which was never trained on inputs with
  a $\mathbf{1}$ component, and a hazard of the raw convention. It is not a change of the
  model's features.
\end{itemize}

\begin{table}[t]
  \caption{Contract execution: canonical TransformerLens/SAELens path versus the fast path
  (16 calibration windows; 96 admissible edits for K7). Tolerances were fixed before the
  float32 run; the float64 run is a diagnostic with the same tolerances.}
  \label{tab:contract-exec}
  \centering
  \scriptsize
  \setlength{\tabcolsep}{2.5pt}
  \begin{tabular}{llccc}
\toprule
Check & Quantity & Tol. & float32 (v1) & float64 (v2) \\
\midrule
K0 weights/config identity & bitwise & -- & \checkmark & \checkmark \\
K1 tokenizer (16 windows) & equal & -- & 16/16 & 16/16 \\
K2 activation, $x=h-\bar h$ & rel.\ $\ell_2$ & $10^{-4}$ & $4.2\times10^{-6}$ & $1.0\times10^{-14}$ \\
\quad uncentred $h$ (neg.\ control) & rel.\ $\ell_2$ & -- & 0.0319 & 0.0319 \\
K3 SAE features & rel.\ $\ell_2$ & $10^{-3}$ & $7.7\times10^{-6}$ & $1.6\times10^{-14}$ \\
K4 L0 / FVE (canonical) & band & [30,120] / 0.7 & 72.4 / 0.827 & 72.4 / 0.827 \\
\quad uncentred input & -- & -- & 176 / 0.003 & 176 / 0.003 \\
K5 no-op hooks & max $|\Delta|$ & 0 & 0 & 0 \\
K6 common logit shift & max $|c_t|$ & -- & 279 & 279 \\
\quad residual after shift & max & $10^{-3}$ & $4.3\times10^{-4}$ & $7.4\times10^{-13}$ \\
\quad probabilities & max $|\Delta p|$ & $10^{-5}$ & $2.8\times10^{-5}$ & $5.1\times10^{-14}$ \\
\quad KL & max & $10^{-6}$ & $7.5\times10^{-7}$ & $1.5\times10^{-15}$ \\
K7 post-edit feature change & rel.\ $\ell_2$ & $10^{-3}$ & $7.6\times10^{-5}$ & $1.9\times10^{-13}$ \\
\quad post-edit probabilities & max $|\Delta p|$ & $10^{-5}$ & $2.8\times10^{-5}$ & $5.3\times10^{-14}$ \\
\quad post-edit KL & max & $10^{-6}$ & $7.7\times10^{-7}$ & $1.5\times10^{-15}$ \\
K9 mean-only edit, model KL & max & $10^{-8}$ & $6.0\times10^{-7}$ & $1.3\times10^{-15}$ \\
\quad SAE drift / $\rho$, raw off-slice $E(x{+}\delta)$ & median & -- & 125 & 125 \\
\quad SAE drift / $\rho$, canonical $E(P(h{+}\delta))$ & max & -- & not run & $3.8\times10^{-15}$ \\
K10 no-processing path, KL & max & $10^{-6}$ & not run & $1.9\times10^{-15}$ \\
\midrule
Verdict & & & \textsc{fail} (blocked) & \textsc{pass} (float64) \\
\bottomrule
\end{tabular}

\end{table}

%%%%%%%%%%%%%%%%%%%%%%%%%%%%%%%%%%%%%%%%%%%%%%%%%%%%%%%%%%%%%%%%%%%%%%%%%%%%%%
\section{Admissible Feature Realization}
\label{sec:fidelity}

\subsection{What linear algebra already fixes}
\label{sec:analysis}
Several facts are standard or elementary, and they decide which outcomes are forced.
Formal statements are in \cref{app:facts} and proofs in \cref{app:proofs}. In the real
model, $E$ stands for $EP$, the Jacobian of the canonical readout (\cref{sec:coord}).
\begin{itemize}
  \item \emph{Rescaling} fixes the target coordinate. It removes leakage on $Q$ only if
  column $j$ of $ED$ restricted to $S\cup Q$ is a multiple of $e_j$ (\cref{prop:mismatch}).
  \item \emph{A least-norm correction} $M^{+}t$ exists exactly when $t\perp\ker M^{\top}$.
  Otherwise the residual certifies infeasibility; generic targets are infeasible once
  $|S\cup Q|>\operatorname{rank}M$ (\cref{prop:feasible,cor:regimes}).
  \item \emph{For ReLU}, with protected pre-activations held fixed, all remaining drift is
  newly activated drift from inactive features that cross zero. A post-activation Jacobian
  cannot activate an inactive target (\cref{prop:relu}).
  \item \emph{JumpReLU and TopK} add range gaps and evictions. Their local derivatives are
  not global inverses (\cref{prop:jump}).
  \item \emph{Norm budgets} give a trust-region problem with a ridge solution
  (\cref{prop:tr}).
  \item \emph{Cost.} For same-layer edits the constraint matrix is just the gathered
  encoder rows. Matrix-free solvers matter only for downstream readouts
  \citep{paige1982lsqr}.
\end{itemize}

\subsection{Real-model pilot}
\label{sec:pilot}
\paragraph{Protocol (P3-REAL-01R pilot).} Fixed before any test document was read:
calibration and test sets are WikiText-103 validation and test articles; 64 features are drawn
from a firing-density band on calibration windows, and the pilot uses a seeded subsample of
\PilotFeatures; per-feature target changes are the calibration activation quantiles
$q_{50}, q_{90}, q_{99}$ and $2q_{99}$, and equal-norm budgets are fractions of the median
residual norm, all frozen with a hash. Each feature has up to 8 test documents with an
active-target point and 8 with an inactive-target point. All edits are admissible and are
compared at equal norm or at matched target change; unmatched rows are kept. Quality is
KL$(P_{\mathrm{clean}}\,\|\,P_{\mathrm{edit}})$ and $\Delta$NLL over the 16 positions after
the edit, for 2 documents per feature and target type. The unit is the test document, with
95\% document-cluster bootstrap intervals. The run (float64, two CPU threads) produced
\PilotRows{} rows from \PilotDocs{} documents; the median \LN{} solve took
\PilotSolveMs\,ms.

\paragraph{Results.}
\Cref{tab:pilot} reports the $q_{99}$ dose.
\begin{itemize}
  \item \emph{Decoder.} For active targets, the target-matched decoder edit leaves
  $D_{\mathrm{orig}} = \PilotOrigDec\alpha$ and $D_{\mathrm{all}} = \PilotAllDec\alpha$.
  The weight-only rescaled decoder is identical here, because pre-activations are affine.
  For inactive targets it misses the target (error \PilotRescInactTgt).
  \item \emph{\LN{} correction.} It removes $D_{\mathrm{orig}}$ by construction. What
  remains is newly activated drift, $D_{\mathrm{all}} = \PilotAllLn\alpha$ against
  \PilotAllDec$\alpha$ for the decoder.
  \item \emph{Per-document paired differences} (\LN{} minus decoder, active targets):
    \begin{itemize}
      \item $\Delta D_{\mathrm{all}} = \PilotDAllAct$;
      \item $\Delta D_{\mathrm{new}} = \PilotDNewAct$;
      \item $\Delta\mathrm{KL} = \PilotDKLAct$ nats.
    \end{itemize}
    For active targets the same signs hold at every dose (\cref{tab:pilot-paired}). For
    inactive targets, $\Delta D_{\mathrm{all}}$ and $\Delta$KL are negative at every dose,
    but the $\Delta D_{\mathrm{new}}$ intervals include zero at $q_{50}$ and $q_{90}$.
  \item \emph{Norm and encoder row.} The \LN{} edit has a smaller norm than the decoder edit
  (\PilotNormLn{} against \PilotNormDec) and is only slightly larger than the encoder-row
  edit (difference \PilotLnEncNorm). Against the encoder row, it lowers
  $D_{\mathrm{all}}$ by \PilotLnEncAll, with a KL difference of \PilotLnEncKL.
  \item \emph{Controls.} Random directions could not reach the target in \PilotRandInf{} of
  \PilotRandTotal{} target-matched rows, which are kept as infeasible (\cref{sec:denominators}).
  The stored mean-only rows use the raw off-slice readout. Canonically a mean-only edit
  changes no feature (\cref{sec:coord}), so it has no target-matched row and is excluded
  from every table and ranking.
\end{itemize}
\Cref{fig:pilot-dose} shows the dose dependence. These are internal measurements, and
\cref{sec:external} explains why they do not show external benefit.

\begin{table*}[t]
  \caption{Real-model pilot (GPT-2 small, residual SAE at layer 8, float64). Target-matched
  edits at dose $q_{99}$; medians over test documents. Drift is relative to $\alpha$. KL is
  over the 16 positions after the edit (2 documents per feature and target type), with
  $n_{\mathrm{doc}}$ given for drift. All target-matched methods reach the target (error
  below $10^{-7}$) except the weight-only rescaled decoder for inactive targets. Mean-only
  rows are omitted (raw off-slice readout; \cref{sec:coord}). Document-cluster 95\%
  intervals are in \cref{tab:pilot-ci}.}
  \label{tab:pilot}
  \centering
  \small
  \begin{tabular}{llccccccc}
\toprule
Target & Method & matched & $n_{\mathrm{doc}}$ & $D_{\mathrm{orig}}/\alpha$ & $D_{\mathrm{new}}/\alpha$ & $D_{\mathrm{all}}/\alpha$ & KL$\times10^{3}$ & $\|\delta\|$ \\
\midrule
active & decoder & 64/64 & 42 & 0.583 & 0.311 & 0.656 & 5.32 & 21.6 \\
 & rescaled (weight-only) & 64/64 & 42 & 0.583 & 0.311 & 0.656 & 5.32 & 21.6 \\
 & encoder row & 64/64 & 42 & 0.681 & 0.366 & 0.778 & 1.03 & 11.8 \\
 & LN ($J_EP$) & 64/64 & 42 & $\approx 0$ & 0.114 & 0.114 & 0.96 & 12.7 \\
 & random & 31/64 & 24 & 7.33 & 108 & 108 & 691 & 438 \\
\midrule
inactive & decoder & 64/64 & 40 & 0.583 & 0.361 & 0.716 & 3.41 & 36.9 \\
 & rescaled (weight-only) & 0/64 & 40 & 0.342 & 0.112 & 0.37 & 0.65 & 22.5 \\
 & encoder row & 64/64 & 40 & 0.943 & 0.943 & 1.36 & 1.38 & 21 \\
 & LN ($J_EP$) & 64/64 & 40 & $\approx 0$ & 0.295 & 0.295 & 1.16 & 22.5 \\
 & random & 29/64 & 23 & 17.4 & 347 & 348 & 327 & 1174 \\
\bottomrule
\end{tabular}

\end{table*}

\begin{figure}[t]
  \centering
  \begin{tikzpicture}
    \begin{groupplot}[
        group style={group size=2 by 1, horizontal sep=0.9cm},
        width=0.58\columnwidth, height=3.6cm,
        xtick={0,1,2,3}, xticklabels={$q_{50}$,$q_{90}$,$q_{99}$,$2q_{99}$},
        ymode=log, grid=major, tick label style={font=\scriptsize},
        label style={font=\scriptsize}, title style={font=\scriptsize}]
      \nextgroupplot[title={active target}, ylabel={median drift / $\alpha$},
        legend to name=pilotleg, legend columns=3, legend style={font=\scriptsize}]
      \addplot[black, thick, mark=o] table[x=dose_idx, y=active_decoder_all] {figures/pilot_dose.dat};
      \addlegendentry{decoder: $D_{\mathrm{all}}$}
      \addplot[black, densely dotted, mark=o] table[x=dose_idx, y=active_decoder_new] {figures/pilot_dose.dat};
      \addlegendentry{decoder: $D_{\mathrm{new}}$}
      \addplot[blue, thick, mark=square*] table[x=dose_idx, y=active_jacobian_ln_all] {figures/pilot_dose.dat};
      \addlegendentry{\LN: $D_{\mathrm{all}}{=}D_{\mathrm{new}}$}
      \nextgroupplot[title={inactive target}]
      \addplot[black, thick, mark=o] table[x=dose_idx, y=inactive_decoder_all] {figures/pilot_dose.dat};
      \addplot[black, densely dotted, mark=o] table[x=dose_idx, y=inactive_decoder_new] {figures/pilot_dose.dat};
      \addplot[blue, thick, mark=square*] table[x=dose_idx, y=inactive_jacobian_ln_all] {figures/pilot_dose.dat};
    \end{groupplot}
  \end{tikzpicture}
  \\ \ref{pilotleg}
  \caption{Pilot: drift versus calibration dose (target-matched; medians over test
  documents). The \LN{} correction has no originally-active drift, so all of its drift is
  newly activated.}
  \label{fig:pilot-dose}
\end{figure}

\subsection{When the correction redistributes drift}
\label{sec:redistribution}
A correction that protects some features can only move drift elsewhere: onto unprotected
features in a linear model, or onto features that cross threshold in a ReLU model
(\cref{prop:relu}). Whether it also \emph{reduces} total drift depends on the regime, which
controlled toy studies isolate. These results are unfavorable to the correction in the
dense regime and are reported here, not in the appendix.

\paragraph{Controlled toy study: design.}
\label{sec:toy}
The toy study (Python standard library) fixed its hypotheses, metrics, unit (one random draw
of $(D, E, b)$), budget and stop conditions in a configuration before the reported run;
development seeds 0--2 were used only for debugging, and seeds 1000--1019 were run once (20
instances per cell). One change followed the development seeds: an exploratory sparsity
axis (\cref{app:toy}). All \FeasTotalPass/\FeasTotal{} feasibility checks and verification
identities passed (\cref{tab:feas}); they show that the code is correct and are not findings.
Every number below is re-aggregated from the raw files.

\paragraph{Linear regimes (toy): rescaling versus interference.}
\label{sec:linear}
\Cref{tab:linear} (\cref{app:tables}) places seven linear regimes into the classes predicted by
\cref{prop:mismatch,prop:feasible}. A class is assigned only if it holds in all 20
instances.
\begin{itemize}
  \item \emph{R0} has no mismatch, yet the encoder-row edit leaks \EncGradLeakRzero.
  \emph{R1} is a pure scale error (decoder target error \RoneDecTgt), which rescaling fixes.
  \item \emph{R2} is pure interference. The leakage is \RtwoDecLeak{} $= 0.3\sqrt7$, and
  rescaling cannot change it. The \LN{} correction removes it at norm ratio \RtwoLnNorm{}
  $=\sqrt{(G^{-1})_{jj}}$. At the decoder's norm, the trust region reaches target error
  \RtwoTrTgt{} with leakage \RtwoTrLeak.
  \item \emph{R4a} protects all 47 non-targets with $d=16$ and is infeasible (median
  relative residual \RfouraRes; every row at least \RfouraResRowMin). The least-squares
  solution gives up the target (error \RfouraLnTgt).
  \item \emph{R4b and R5} protect 7 features. This zeroes their drift but \emph{raises}
  total leakage. In R4b it goes from \RfourbDecLeak{} to \RfourbLnLeak{} at norm ratio
  \RfourbLnNorm. In R5 it goes from \RfiveRescLeak{} (rescaled) to \RfiveLnLeak{} (\LN).
\end{itemize}


\paragraph{Active sets and norm budgets (toy, pre-registered).}
\label{sec:activeset}
This setting has $d=16$ and $m=64$, an encoder halfway between tied and biorthogonal, and
base points built from three decoder atoms. The median number of active features is
\DenseReluActMed{} (range \DenseReluActMin--\DenseReluActMax) for ReLU and
\DenseJumpActMed{} for JumpReLU. This is dense: $|S\cup Q|$ is a large fraction of $d$.
\Cref{tab:relu} (\cref{app:tables}) gives the ReLU decomposition.

The \LN{} correction removes most originally active drift. Newly activated drift grows,
however, because the correction's larger norm pushes inactive features over threshold.
Its total drift was lower than the rescaled decoder's in \DenseReluActQWins,
\DenseReluActOneWins{} and \DenseReluActFourWins{} instances for active targets at
$\alpha = 0.25, 1, 4$, and in \DenseReluInactQWins, \DenseReluInactOneWins{} and
\DenseReluInactFourWins{} for inactive targets.

Bounded repair reached a crossing-free edit in \RepairConvRJMin--\RepairConvRJMax{} of
\RepairRowsPerCell{} rows per ReLU or JumpReLU cell, but it did so by abandoning the target:
median target error \RepairTgtMin--\RepairTgtMax{} over range-feasible cells. For TopK it
converged in only \RepairConvTKMin--\RepairConvTKMax{} of \RepairRowsPerCell{} rows per cell,
because eviction is rank-based (\cref{prop:jump}). As \cref{prop:relu} predicts, the naive
post-activation Jacobian never moved an inactive target. With JumpReLU, an inactive target
at $\alpha = 0.25<\theta = 0.3$ was unattainable for every method.

At the decoder's norm, the trust region had lower combined error than the decoder in
\TrWinsMin--\TrWinsMax{} of 20 instances per ReLU/JumpReLU cell (\cref{tab:budget}). This
comparison is circular: the trust region minimizes the same internal quantity it is scored
on.


\paragraph{Sparser active sets (toy, exploratory).}
\label{sec:sparse}
This subsection is exploratory. The axis was added after inspecting development seeds,
so it should not be read as confirmation.

Changing only the ReLU bias lowers the median active count to \SparseFiftyActMed{}
(bias $-0.5$) and \SparseSeventyFiveActMed{} (bias $-0.75$). \Cref{fig:alpha} and
\cref{tab:sparse} show the effect.
\begin{itemize}
  \item With bias $-0.75$, the \LN{} correction had lower total drift than rescaling in
  \SparseSeventyFiveActQWins, \SparseSeventyFiveActOneWins{} and
  \SparseSeventyFiveActFourWins{} instances for active targets.
  \item With bias $-0.5$, it was lower in \SparseFiftyActQWins{} and \SparseFiftyActOneWins{}
  instances at $\alpha\le1$, but only \SparseFiftyActFourWins{} at $\alpha = 4$.
\end{itemize}
In every panel of \cref{fig:alpha}, the drift that remains after the \LN{} correction is
almost entirely newly activated drift.

Our reading, which is a hypothesis about real SAEs and not a finding: a local correction
helps only when the protected set is small relative to $d$ \emph{and} the edit is small
enough that few inactive features cross threshold.

\begin{figure*}[t]
  \centering
  \begin{tikzpicture}
    \begin{groupplot}[
        group style={group size=3 by 1, horizontal sep=1.4cm},
        width=0.33\textwidth, height=4.0cm,
        xmode=log, log basis x=2, xtick={0.25,1,4}, xticklabels={0.25,1,4},
        xlabel={$\alpha$}, ymin=0, grid=major, tick label style={font=\scriptsize},
        label style={font=\small}, title style={font=\small}]
      \nextgroupplot[title={dense (pre-registered)}, ylabel={median drift / $\alpha$},
        legend to name=alphaleg, legend columns=4, legend style={font=\scriptsize}]
      \addplot[black, thick, mark=o] table[x=alpha, y=dense_resc_total] {figures/alpha_curves.dat};
      \addlegendentry{rescaled: total}
      \addplot[black, densely dotted, mark=o] table[x=alpha, y=dense_resc_orig] {figures/alpha_curves.dat};
      \addlegendentry{rescaled: orig.\ active}
      \addplot[blue, thick, mark=square*] table[x=alpha, y=dense_ln_total] {figures/alpha_curves.dat};
      \addlegendentry{\LN: total}
      \addplot[blue, dashed, mark=square*] table[x=alpha, y=dense_ln_new] {figures/alpha_curves.dat};
      \addlegendentry{\LN: newly active}
      \nextgroupplot[title={bias $-0.5$ (exploratory)}]
      \addplot[black, thick, mark=o] table[x=alpha, y=bm050_resc_total] {figures/alpha_curves.dat};
      \addplot[black, densely dotted, mark=o] table[x=alpha, y=bm050_resc_orig] {figures/alpha_curves.dat};
      \addplot[blue, thick, mark=square*] table[x=alpha, y=bm050_ln_total] {figures/alpha_curves.dat};
      \addplot[blue, dashed, mark=square*] table[x=alpha, y=bm050_ln_new] {figures/alpha_curves.dat};
      \nextgroupplot[title={bias $-0.75$ (exploratory)}]
      \addplot[black, thick, mark=o] table[x=alpha, y=bm075_resc_total] {figures/alpha_curves.dat};
      \addplot[black, densely dotted, mark=o] table[x=alpha, y=bm075_resc_orig] {figures/alpha_curves.dat};
      \addplot[blue, thick, mark=square*] table[x=alpha, y=bm075_ln_total] {figures/alpha_curves.dat};
      \addplot[blue, dashed, mark=square*] table[x=alpha, y=bm075_ln_new] {figures/alpha_curves.dat};
    \end{groupplot}
  \end{tikzpicture}
  \\ \ref{alphaleg}
  \caption{Where drift goes (ReLU, active targets, 20 instances, medians). In the dense
  setting the \LN{} correction trades originally active drift for more newly activated
  drift. In the sparser exploratory settings it lowers total drift at small $\alpha$.}
  \label{fig:alpha}
\end{figure*}


\paragraph{Where the real SAE falls.}
\label{sec:redist-real}
In the pilot, the SAE has about \CxLZero{} active features per token out of $d=768$, a
protected fraction far below that of the dense toy. Its behavior resembles the sparse
exploratory settings rather than the dense pre-registered one.
\begin{itemize}
  \item The \LN{} correction's newly activated drift did not exceed the decoder's: for
  active targets $\Delta D_{\mathrm{new}} = \PilotDNewAct$ at $q_{99}$
  (\cref{tab:pilot-paired}).
  \item Its total drift was lower at every dose.
\end{itemize}
This fits the regime hypothesis stated after the exploratory toy study. It is a pilot with
\PilotFeatures{} features and one layer, however, so it is consistent with that hypothesis
rather than a test of it. The full calibration set of 64 features and other layers or SAE
families are needed before any regime claim. \todo{P3-REAL-01R-FULL: all 64 calibration
features, same frozen doses, same document-paired analysis}.

%%%%%%%%%%%%%%%%%%%%%%%%%%%%%%%%%%%%%%%%%%%%%%%%%%%%%%%%%%%%%%%%%%%%%%%%%%%%%%
\section{Matched-Denominator Robustness}
\label{sec:denominators}
A target-matched comparison is only as good as its denominator. We re-aggregated the stored
pilot rows (no rerun) with rules fixed in the aggregation script before its tables were
produced: a row is matched if its status is OK and its target error is at most $10^{-6}$;
unmatched and infeasible rows are counted, never dropped or pooled
(\cref{tab:denominators}).
\begin{itemize}
  \item \emph{Primary methods share one denominator.} Decoder, encoder row and \LN{} are
  matched at every test point and dose (largest target error \RgTgtErrMax, float32 rounding of
  the stored reference), so the comparisons of \cref{sec:pilot} use all test points
  (exclusion \RgPrimExclMax).
  \item \emph{Random directions lose half.} No non-negative scale reaches the target for
  \RgSecExclMin--\RgSecExclMax{} of test points per dose and target type. Adding random to
  the intersection would halve the sample, so it is summarized only on its own feasible rows
  and never ranked against the others.
  \item \emph{Unmatched baselines stay separate.} The weight-only rescaled decoder lands on
  the target in all \RgRescActWithin{} active-target rows, because pre-activations are affine,
  and in \RgRescInactWithin{} inactive-target rows (median error \RgRescInactErr{} at
  $q_{99}$). The plain decoder $\alpha D_j$ is never matched.
  \item \emph{Mean-only rows exist only under the raw readout.} All \RgRawOnly{} of them
  are excluded.
  \item \emph{One earlier readout mixed denominators.} Gate G1 (rescaled originally-active
  drift at $q_{99}$) pooled both target types: \RgGOneAsRun{} over 54 documents. Split, it
  is \RgGOneAct{} for active targets (matched, $n=\RgGOneActN$) and \RgGOneInact{} for
  inactive targets (unmatched, $n=\RgGOneInactN$). Both lie above $0.1\alpha$, so the
  readout's side is unchanged, but the pooled number is not a matched quantity.
  \item \emph{In-run amendments.} The rule that a gate readout needs at least 10 documents
  and the derived target-gain metric were committed while the pilot's test run was
  executing, after its configuration was frozen and before any test row existed on disk. We
  label both as in-run amendments. Target gain is no longer used in the body: its headline
  value came from the raw mean-only readout.
  \item \emph{Effective calibration counts.} Each pilot feature's $q_{99}$ is a quantile of
  \RgNPosMin--\RgNPosMax{} positive calibration activations. For \RgTopTwo{} of the
  \PilotFeatures{} features there are at most 101, so $q_{99}$ lies between the two largest
  values and is effectively a top order statistic. Of all 64 calibration features,
  \RgAllBelowHundred{} have fewer than 100 positives (minimum \RgAllNPosMin). We keep the
  frozen $q_{99}$ and do not switch doses after the fact.
\end{itemize}

\begin{table}[t]
  \caption{Pilot row status by method (target-matched mode, all four doses, active and
  inactive targets). ``Raw-only'': rows that exist only under the raw off-slice readout.
  ``Docs'': largest number of test documents in one dose and target-type cell.}
  \label{tab:denominators}
  \centering
  \scriptsize
  \setlength{\tabcolsep}{3pt}
  \begin{tabular}{lrrrrrr}
\toprule
Method (target-matched, 4 doses) & rows & matched & unmatched & infeasible & raw-only & docs \\
\midrule
decoder & 512 & 512 & 0 & 0 & 0 & 42 \\
encoder row & 512 & 512 & 0 & 0 & 0 & 42 \\
LN ($J_EP$) & 512 & 512 & 0 & 0 & 0 & 42 \\
random & 512 & 240 & 0 & 272 & 0 & 42 \\
rescaled (weight-only) & 512 & 0 & 512 & 0 & 0 & 42 \\
plain decoder ($\alpha D_j$) & 512 & 0 & 512 & 0 & 0 & 42 \\
mean-only (raw readout) & 512 & 0 & 0 & 0 & 512 & 42 \\
\bottomrule
\end{tabular}

\end{table}

%%%%%%%%%%%%%%%%%%%%%%%%%%%%%%%%%%%%%%%%%%%%%%%%%%%%%%%%%%%%%%%%%%%%%%%%%%%%%%
\section{Lexical Change and Quality}
\label{sec:external}

Everything above is measured with the encoder that defines the edit. This section states
what that can certify and then reports one limited behavioral pilot.

\subsection{What agreement with the encoder can certify}
\label{sec:certify}
Where it is feasible, the \LN{} correction makes $D_{\mathrm{orig}} = 0$ and the target error
zero by construction, as measured with the same encoder. That says nothing about features
outside $S\cup Q$ (\cref{prop:relu}) or about any other readout. A synthetic example makes
this concrete. A ground-truth dictionary $A$ ($d=16$, $m=8$, pairwise cosine 0.3) defines an
external readout $z = A^{+}x$, and the SAE is a noisy estimate of $A$
(\cref{tab:proxy}, \cref{app:tables}). The \LN{} correction drives the internal error below
$10^{-6}$ in every case, yet the external error falls from \TxEncAccDecExt{} to
\TxEncAccLnExt{} when the encoder is exact (better in \TxEncAccWins{} instances), rises from
\TxDecAccDecExt{} to \TxDecAccLnExt{} when the decoder is exact (better in \TxDecAccWins),
and is unchanged when both are noisy (better in \TxBothWins). These outcomes are fixed by
construction: the internal metric measures agreement with the encoder. Behavior can also
change while SAE readouts are held fixed \citep{cui2026sae}.

\subsection{Lexical pilot: blocked at calibration}
\label{sec:real02}
We froze one limited behavioral test (P3-REAL-02-LX, \cref{app:real02}) before calibration.
\begin{itemize}
  \item \emph{Design.} At most 32 held-out WikiText-103 articles, one keyword-free prompt
  each, greedy decoding of 32 tokens, and the edit at every position. The methods are the
  projected decoder, the projected encoder row, the \LN{} correction ($J_EP$) and projected
  DiffMean, each at the largest norm budget whose calibration KL on the common prompt is at
  most $0.1$ nats, plus no edit. The mean-only control runs at one fixed amplitude outside
  every ranking.
  \item \emph{Metric.} \emph{Wedding-lexicon occurrence}: whether the continuation contains
  an ActAdd keyword \citep{turner2024steering}, and hits per generated token. This is a
  lexical proxy. The lexicon also selects the feature and the DiffMean direction, so the
  metric is not independent of direction selection.
\end{itemize}
\emph{Outcome.} Choosing the feature and DiffMean needs at least 20 calibration windows that
contain a keyword. With the frozen cap of 8 windows per article, \LxHitsVOne{} qualified and
the stage returned \textsc{blocked}. We kept that record. Before reading any test document,
we committed one amendment: scan every calibration window, with a second
\textsc{blocked} ending the experiment. All \LxWindowsAll{} windows then gave
\LxHitsVOneOne. P3-REAL-02-LX is therefore \textsc{blocked}. No feature was chosen, no test
document was read, no text was generated, and no lexical, KL or NLL number exists. We did not
lower the minimum after seeing the count.
Unblocking it needs a new frozen protocol: either calibration text with more topic support,
which means new data and a download approval, or a topic that WikiText supports. Separately,
the adapter was checked on a synthetic prompt. Greedy decoding matches the library's
\texttt{generate}, cached \LN{} directions equal recomputation, actor edits have mean
component below $10^{-12}$, and the mean-only control leaves the continuation unchanged.
\todo{P3-REAL-02-LX: held-out wedding-lexicon occurrence and its KL/NLL tradeoff; blocked
at calibration, needs a new frozen protocol}
\todo{P3-REAL-02-SEM: an independent semantic evaluation of the same continuations}.

This paper therefore contains no behavioral evidence for or against the correction.

%%%%%%%%%%%%%%%%%%%%%%%%%%%%%%%%%%%%%%%%%%%%%%%%%%%%%%%%%%%%%%%%%%%%%%%%%%%%%%
\section{Limitations}
\label{sec:limitations}
\begin{itemize}
  \item \emph{One model, one SAE, one layer.} The fidelity results come from a pilot with
  \PilotFeatures{} features on \PilotDocs{} documents and are conditional on that sample.
  WikiText differs from the SAE's OpenWebText training data, and quality metrics come from 2
  documents per feature and target type.
  \item \emph{The invariance proof is GPT-2 specific.} Architectures whose residual readers
  are not mean-subtracting need their own analysis.
  \item \emph{float32 is characterized, not fixed.} The float32 contract failed its
  pre-registered tolerance. We ran everything in float64.
  \item \emph{One unrecorded detail.} The TransformerLens version used to train the SAE is
  not recorded; our reconstruction check (L0 \CxLZero, FVE \CxFVE{} against the release's 60
  and 0.9) supports the coordinates but cannot pin down the training code.
  \item \emph{Toy results are small and largely forced}; the sparse axis is exploratory.
  \item \emph{Internal numbers are internal.} KL and $\Delta$NLL measure total output change,
  not success.
  \item \emph{The behavioral metric is lexical,} and the lexicon also chose the steered
  feature and the DiffMean direction. Semantic success is \notrun.
  \item \emph{Other choices.} The Euclidean norm is a modeling choice
  \citep{park2024linear}. JumpReLU and TopK families \citep{lieberum2024gemma} add the
  boundary effects of \cref{prop:jump} and are not covered.
\end{itemize}

\section{Conclusion}
\label{sec:conclusion}
Auditing an SAE edit starts with the readout. For GPT-2 small and the layer-8 residual SAE,
the SAE's inputs live in $\operatorname{range}(P)$. Reading an edit raw, as $E(x+\delta)$,
can report large feature changes for edits the model cannot see; the canonical readout
$E(P(h+\delta))$ removes them. With that fixed, and with edits restricted to model-visible
directions, a conditional same-layer pilot shows the following:
\begin{itemize}
  \item projected decoder edits leave substantial drift on other features;
  \item a least-norm correction with $J_E(Ph)P$ lowers that drift on a common denominator of
  test points, whereas in a dense toy regime it moves drift onto newly activated features.
\end{itemize}
None of this is evidence of external control. The one behavioral test we froze, a lexical
proxy, was blocked at calibration for lack of topic text, so whether the correction changes
behavior, lexically or semantically, remains open.

\section*{Impact Statement}
This work concerns the reliability of interpretability-based interventions, which are
sometimes used to support safety claims such as removed capabilities or suppressed
behaviors. Checking the readout convention and separating internal from external effects
should reduce over-claiming. A raw readout can report large SAE changes for edits the model
cannot see, which matters for SAE-based monitoring; null-space machinery that holds SAE
readouts fixed also bears on SAE-based defenses \citep{cui2026sae}. We do not develop
attacks, and our behavioral task is benign topic steering.

\bibliography{references}
\bibliographystyle{icml2026}

%%%%%%%%%%%%%%%%%%%%%%%%%%%%%%%%%%%%%%%%%%%%%%%%%%%%%%%%%%%%%%%%%%%%%%%%%%%%%%
\newpage
\appendix
\onecolumn

\section{Standard Facts}
\label{app:facts}
These statements are standard or elementary; none is claimed as new.

\paragraph{Linear encoders.}
\begin{proposition}[Mismatch and rescaling]
\label{prop:mismatch}
If $\sigma$ is the identity, then $\Delta a = \alpha EDe_j$ for $\delta = \alpha D_j$.
Rescaling hits the target if and only if $(ED)_{jj}\neq 0$, and then leaves leakage
$\alpha\|(ED)_{Qj}\|/|(ED)_{jj}|$ on $Q$. It therefore suffices if and only if column $j$ of
$ED$, restricted to $S\cup Q$, is a multiple of $e_j$. For a tied encoder $E = D^{\top}$
whose decoder columns have pairwise cosine $c$, the leakage is $c\sqrt{m-1}\,\alpha$ both
before and after rescaling.
\end{proposition}

\begin{proposition}[Least-norm solutions and infeasibility]
\label{prop:feasible}
$M\delta = t$ is solvable if and only if $t\perp\ker M^{\top}$. If $M$ has full row rank,
$\delta^\star = M^{\top}(MM^{\top})^{-1}t$ is the unique minimum-norm solution, and every
other solution has squared norm $\|\delta^\star\|^2 + \|n\|^2$ for some $n\in\ker M$.
Otherwise the residual $r = t - MM^{+}t$ certifies infeasibility: $M^{\top}r = 0$ and
$r^{\top}t = \|r\|^2 > 0$. In particular, a generic target is infeasible once
$|S\cup Q| > \operatorname{rank}M$, which holds whenever $|S\cup Q| > d$. Two identical
encoder rows with targets $(\alpha,0)$ leave residual $\alpha/\sqrt2$ in any dimension.
\end{proposition}
These are textbook facts \citep{golub2013matrix,bjorck1996numerical}.

\begin{corollary}[Three regimes]
\label{cor:regimes}
Let $Q$ contain all non-target features.
\begin{itemize}
  \item If $E = D^{+}$ and $D$ has full column rank, then $\delta^\star = \alpha D_j$: the
  decoder edit is already least-norm.
  \item If $D$ has orthonormal columns and $E = \operatorname{diag}(s)D^{\top}$, then
  $\delta^\star$ is the rescaled decoder edit.
  \item If $E = D^{\top}$ with Gram matrix $G = D^{\top}D$, then
  $\|\delta^\star\|/\alpha = \sqrt{(G^{-1})_{jj}}\ge 1$, with equality if and only if
  $Ge_j\propto e_j$.
\end{itemize}
\end{corollary}

\begin{remark}[Compute]
For same-layer edits, the Jacobian of $p$ restricted to $S\cup Q$ is the gathered matrix
$E_{S\cup Q}$. With $|S\cup Q|\ll d$ it is small, so materializing it costs little.
Matrix-free solvers that use only Jacobian--vector and vector--Jacobian products, such as
CGLS/LSQR \citep{paige1982lsqr,bjorck1996numerical}, matter when the protected readouts sit at
a later layer. In the toy study, CGLS matched the dense solution to within \VfiveMaxDiff{}
(relative).
\end{remark}

\paragraph{Nonlinear boundaries.}
\begin{proposition}[ReLU piece exactness]
\label{prop:relu}
For ReLU, suppose no pre-activation, including the target's, changes sign between $x$ and
$x+\delta$. Then $\Delta a = \operatorname{diag}(\mathbb{1}[p>0])E\delta$ exactly. If the
pre-activations of $Q = A_0\setminus\{j\}$ are held fixed, then
$D_{\mathrm{kept}} = D_{\mathrm{deact}} = 0$, and all remaining drift is $D_{\mathrm{new}}$. It
comes from inactive features with $(E\delta)_i > -p_i$. The almost-everywhere Jacobian row
of an inactive target is zero, so a solve on post-activation Jacobians cannot activate it.
\end{proposition}

\begin{proposition}[JumpReLU and TopK]
\label{prop:jump}
A JumpReLU feature cannot take values in $(0,\theta]$, and crossing $\theta$ changes the
activation by at least $\theta$ where the almost-everywhere derivative predicts no change.
Under TopK with exactly $k$ positive active features, an inactive target can enter only by
evicting at least one of them. In both cases the local derivative is not a global inverse.
For TopK, comparing active sets at the two endpoints only lower-bounds the number of
changes along the path.
\end{proposition}

\begin{proposition}[Norm budget]
\label{prop:tr}
$\min_{\|\delta\|\le\tau}\|M\delta - t\|$ is convex. Its solution is $M^{+}t$ if
$\|M^{+}t\|\le\tau$. Otherwise it is $(M^{\top}M+\lambda I)^{-1}M^{\top}t$ with the unique
$\lambda>0$ such that the solution has norm $\tau$ \citep{more1983computing}.
\end{proposition}

\section{Proofs}
\label{app:proofs}
\paragraph{\Cref{prop:mismatch}.}
With $\sigma$ the identity, $p(x+\delta) - p(x) = E\delta$. For $\delta = cD_j$ this is
$cED_j$. Choosing $c = \alpha/(ED)_{jj}$ sets the target coordinate to $\alpha$, and the
$Q$ coordinates are then $\alpha(ED)_{Qj}/(ED)_{jj}$. For $E = D^{\top}$ with unit columns
and pairwise cosine $c$, $(ED)_{jj} = 1$ and $(ED)_{ij} = c$, which gives
$c\sqrt{m-1}\,\alpha$.

\paragraph{\Cref{prop:feasible}.}
The range of $M$ is the orthogonal complement of $\ker M^{\top}$. Under full row rank,
$MM^{\top}$ is invertible and $\delta^\star = M^{\top}(MM^{\top})^{-1}t$ lies in
$\operatorname{range}M^{\top}\perp\ker M$. Any other solution is $\delta^\star + n$ with
$n\in\ker M$, so the Pythagorean identity gives minimality. In general $MM^{+}$ is the
orthogonal projector onto $\operatorname{range}M$. Hence $r = t - MM^{+}t\in\ker M^{\top}$
and $r^{\top}t = r^{\top}(r + MM^{+}t) = \|r\|^2$, which is nonzero iff $t\notin\operatorname{range}M$.
For duplicate rows $e$ with targets $(\alpha,0)$, both residual coordinates depend on the
same scalar $u = e^{\top}\delta$. The minimum of $(u-\alpha)^2 + u^2$ is at $u=\alpha/2$ with
value $\alpha^2/2$. The other rows can be satisfied simultaneously when $e$ and those rows
are linearly independent.

\paragraph{\Cref{cor:regimes}.}
If $D$ has full column rank, then $(D^{+})^{+} = D$. For
$E = \operatorname{diag}(s)D^{\top}$ with orthonormal $D$,
$E^{+} = D\operatorname{diag}(s)^{-1}$. For $E = D^{\top}$,
$(D^{\top})^{+} = DG^{-1}$, and
$\|DG^{-1}e_j\|^2 = e_j^{\top}G^{-1}GG^{-1}e_j = (G^{-1})_{jj}$. By Cauchy--Schwarz in the
$G$-inner product, $1 = (e_j^{\top}e_j)^2 \le (e_j^{\top}Ge_j)(e_j^{\top}G^{-1}e_j)$, and
$G_{jj} = 1$.

\paragraph{\Cref{prop:relu}.}
Along $x + s\delta$ for $s\in[0,1]$, each $p_i(s)$ is affine in $s$. Its sign therefore
changes on the segment iff it differs at the endpoints. If no index changes activity, then
$\sigma$ acts as a fixed diagonal 0/1 map on the segment. Protected pre-activations are held
fixed, so protected activations do not change. An inactive feature changes only if $p_i$
becomes positive. For an inactive target, the $j$-th row of
$\operatorname{diag}(\mathbb{1}[p>0])E$ is zero.

\paragraph{\Cref{prop:jump}.}
JumpReLU outputs $p\,\mathbb{1}[p>\theta]$, whose range is $\{0\}\cup(\theta,\infty)$.
Crossing from $p\le\theta$ to $p>\theta$ changes the output by more than $\theta$. Under
TopK, the active set has at most $k$ elements. If $k$ positive features are active and the
target enters, one must leave. The $k$-th order statistic of affine functions is not
monotone along a segment, so endpoint comparison is only a lower bound.

\paragraph{\Cref{prop:tr}.}
The objective is convex quadratic and the constraint set is convex. The KKT conditions are
$(M^{\top}M+\lambda I)\delta = M^{\top}t$ with $\lambda\ge0$ and
$\lambda(\|\delta\|-\tau) = 0$. Convexity rules out the ``hard case'' of indefinite
trust-region problems \citep{more1983computing}.

\paragraph{\Cref{prop:admissible}.}
Write $\operatorname{LN}(v) = \gamma\odot(v-\bar v\mathbf{1})/\sqrt{s(v)^2+\epsilon} + \beta$
with $\bar v = \mathbf{1}^{\top}v/d$ and $s(v)^2 = \|v - \bar v\mathbf{1}\|^2/d$. For any
scalar $c$, $v + c\mathbf{1} - \overline{(v+c\mathbf{1})}\mathbf{1} = v - \bar v\mathbf{1}$,
so the centred vector, $s$, and hence $\operatorname{LN}$ are unchanged. Let $r^{(k)}_t$ be
the residual entering block $k$ at position $t$, and let the edited run add $c_t\mathbf{1}$
to $r^{(\ell)}_t$ for $t$ in some set $T$ ($c_t = 0$ elsewhere). We show by induction on
$k\ge\ell$ that $\tilde r^{(k)}_t = r^{(k)}_t + c_t\mathbf{1}$ for every $t$. The base case is
the edit. In GPT-2, block $k$ computes
$u_t = r_t + \operatorname{Attn}(\operatorname{LN}_1(r_1),\dots,\operatorname{LN}_1(r_t))_t$
(causal attention with learned absolute positions, which enter only at the embedding) and
$r^{(k+1)}_t = u_t + \operatorname{MLP}(\operatorname{LN}_2(u_t))$. By the hypothesis every
$\operatorname{LN}_1$ input differs by a multiple of $\mathbf{1}$, so all queries, keys,
values and attention outputs are unchanged, and $\tilde u_t = u_t + c_t\mathbf{1}$. The same
argument applied to $\operatorname{LN}_2$ gives the claim for $k+1$. After the last block,
$\operatorname{ln\_f}$ removes $c_t\mathbf{1}$, and the logits $W_U\operatorname{ln\_f}(\cdot)$
are unchanged. Dropout is inactive in evaluation mode. The argument uses exactly the
listed structure: a residual reader that is not mean-subtracting (RMSNorm, a direct read of
the residual, or a post-LN block) breaks it, so we make no claim beyond GPT-2. In floating
point the statement holds up to rounding (\cref{tab:contract-exec}, K9).

\emph{Representatives.} The logits are a function of $h$ modulo $\operatorname{span}(\mathbf{1})$
at each edited position, so any linear complement of $\operatorname{span}(\mathbf{1})$
represents model-visible edits. $\operatorname{range}(P)$ is the orthogonal complement and
the one where the SAE's training inputs lie; choosing it is a convention. The canonical
readout $E(P(h+\delta))$ is invariant to $\delta\mapsto\delta+c\mathbf{1}$ because
$P\mathbf{1} = 0$; the raw readout $E(Ph+\delta)$ is not unless $E\mathbf{1} = 0$.

\emph{Effective Jacobian.} With $p(x) = Ex + b$ and $a = \sigma(p)$, the canonical readout of
$\delta$ is $\sigma(EP(h+\delta)+b)$; its pre-activation Jacobian in $\delta$ is $EP$ and its
almost-everywhere post-activation Jacobian is
$\operatorname{diag}(\sigma'(p(Ph)))EP = J_E(Ph)P$. For the \LN{} correction,
$(MP)^{+}t\in\operatorname{range}((MP)^{\top}) = \operatorname{range}(PM^{\top})
\subseteq\operatorname{range}(P)$, and for $\delta\in\operatorname{range}(P)$ we have
$M\delta = MP\delta$, so the constraint is met exactly whenever it is feasible within
$\operatorname{range}(P)$.

\section{Controlled Study Details}
\label{app:toy}
\paragraph{Pre-registration and deviations.}
The configuration fixed the following before the reported run: verification hypotheses
V1--V8, metrics, units, seeds, a 120\,s and two-thread budget, and stop conditions.
V1--V5 and V7 are identities; failing any of them would indicate a bug. V6 and V8 are
descriptive.

After running development seeds 0--2, we made changes before the reported run (v1 to
v2):
\begin{itemize}
  \item the CGLS tolerance of one feasibility check was made explicit ($10^{-6}$);
  \item a low-rank-rows check was added;
  \item the ReLU bias axis of \cref{sec:sparse} was added and marked exploratory.
\end{itemize}
The v1 sections were not modified.

A later bug in the summary code dropped Boolean flags from aggregates and reported a
median-only external-improvement flag. We fixed it and reran. The raw files were
bit-identical to the first run (SHA-256 recorded in the run manifest).

\paragraph{Regimes and methods.}
\begin{itemize}
  \item Linear regimes use $d=16$ with $m=8$ (equiangular, orthonormal) or $m=48$
  (Gaussian columns). Encoders are pseudo-inverse, tied, or scaled tied, with scales
  log-uniform in $[0.5, 2]$.
  \item The active-set setting uses $d=16$, $m=64$ and
  $E = \operatorname{diag}(s)(\tfrac12D^{\top} + \tfrac12D^{+})$ with $s$ in $[0.7,1.4]$.
  Base points are sums of three decoder atoms with coefficients in $[0.5,1.5]$ plus noise
  $0.02$.
  \item Activation settings: ReLU with bias $-0.25$; JumpReLU with $\theta=0.3$; TopK with
  $k=6$.
  \item Doses are $\alpha\in\{0.25,1,4\}$, and budget multipliers are $\{0.5,1,2\}$ times
  $\|\alpha D_j\|$.
\end{itemize}

\paragraph{Compute ledger.}
\begin{itemize}
  \item Development runs on seeds 0--2 took 8.4\,s and 11.3\,s.
  \item A component timing smoke took about 11\,s.
  \item The reported run took 77.1\,s; the rerun after the summary fix took 76.0\,s.
  \item Regenerating the paper's tables from the raw files takes about 1\,s.
\end{itemize}
All of these ran single-threaded on a CPU.

\section{Additional Tables}
\label{app:tables}
\begin{table*}[h]
  \caption{Linear regimes (20 instances each; medians of per-instance means). Leakage is
  $\|\Delta a_{-j}\|/\alpha$. Norms (in parentheses) are relative to $\|\alpha D_j\|$.
  ``Resid.'' is the relative residual of the \LN{} system.}
  \label{tab:linear}
  \centering
  \small
  \begin{tabular}{lcccccc}
\toprule
Regime & Class & Dec.\ tgt err & Dec.\ leak & Resc.\ leak (norm) & LN leak (norm) & LN resid. \\
\midrule
R0: $E{=}D^{+}$, $m{<}d$ & no mismatch & $<\!10^{-6}$ & $<\!10^{-6}$ & $<\!10^{-6}$ (1) & $<\!10^{-6}$ (1) & $<\!10^{-6}$ \\
R1: orth.\ $D$, $E{=}\mathrm{diag}(s)D^{\top}$ & rescale suffices & 0.365 & $<\!10^{-6}$ & $<\!10^{-6}$ (1.13) & $<\!10^{-6}$ (1.13) & $<\!10^{-6}$ \\
R2: tied, $c{=}0.3$ & removed by LN & $<\!10^{-6}$ & 0.794 & 0.794 (1) & $<\!10^{-6}$ (1.14) & $<\!10^{-6}$ \\
R3: scaled tied, $c{=}0.3$ & removed by LN & 0.34 & 0.884 & 0.989 (1.09) & $<\!10^{-6}$ (1.24) & $<\!10^{-6}$ \\
R4a: tied, $m{>}d$, $P$=all & infeasible & $<\!10^{-6}$ & 1.73 & 1.73 (1) & 0.466 (0.398) & 0.821 \\
R4b: tied, $m{>}d$, $|P|{=}7$ & moved to unprot. & $<\!10^{-6}$ & 1.73 & 1.73 (1) & 2.16 (1.35) & $<\!10^{-6}$ \\
R5: $E{=}D^{+}$, $m{>}d$, $|P|{=}7$ & moved to unprot. & 0.669 & 0.467 & 1.44 (3.1) & 2.21 (3.6) & $<\!10^{-6}$ \\
\bottomrule
\end{tabular}

\end{table*}

\begin{table}[h]
  \caption{Feasibility and boundary checks (20 seeds; counts of passing rows).}
  \label{tab:feas}
  \centering
  \small
  \begin{tabular}{lc}
\toprule
Check & Pass \\
\midrule
Full row rank: least-norm, minimality, CGLS = dense & 80/80 \\
More constraints than $d$: residual $>0$ with certificate & 40/40 \\
Rank-deficient, target in range: residual $\approx 0$ & 40/40 \\
Low-rank rows ($k{=}8$, rank 5): certificate & 20/20 \\
Duplicate rows, targets $(\alpha,0)$: residual $\alpha/\sqrt2$ & 20/20 \\
ReLU: negative target unattainable & 20/20 \\
JumpReLU: target in $(0,\theta]$ unattainable & 20/20 \\
TopK: entering target evicts a feature & 20/20 \\
\bottomrule
\end{tabular}

\end{table}

\begin{table}[h]
  \caption{Norm budgets (pre-registered dense ReLU, active target, $\alpha = 1$; medians over
  20 instances). $\tau = \|\alpha D_j\|$.}
  \label{tab:budget}
  \centering
  \small
  \begin{tabular}{lccc}
\toprule
Method & $\|\delta\|/\tau$ & Target err & Total drift \\
\midrule
decoder & 1 & 0.339 & 0.597 \\
rescaled decoder & 1.57 & $<\!10^{-6}$ & 0.944 \\
LN correction & 3.17 & 0.00955 & 1.74 \\
LN, scaled to $\|\alpha D_j\|$ & 1 & 0.629 & 0.278 \\
trust region, $0.5\tau$ & 0.5 & 0.676 & 0.202 \\
trust region, $\tau$ & 1 & 0.4 & 0.361 \\
trust region, $2\tau$ & 1.93 & 0.137 & 0.819 \\
LN + repair & 0.935 & 0.501 & 0.301 \\
\bottomrule
\end{tabular}

\end{table}

\begin{table}[h]
  \caption{Pre-registered active-set setting for all activations (same columns as
  \cref{tab:relu}). $^\dagger$: target value range-infeasible (JumpReLU, $\alpha<\theta$).}
  \label{tab:activeset-all}
  \centering
  \small
  \begin{tabular}{llc|ccc|cccc|cc}
\toprule
 & & & \multicolumn{3}{c|}{Rescaled decoder} & \multicolumn{4}{c|}{LN correction} & \multicolumn{2}{c}{LN + repair} \\
Act. & Target & $\alpha$ & orig.\ act. & new & norm & tgt err & orig.\ act. & new & norm & tgt err & total \\
\midrule
ReLU & active & 0.25 & 0.827 & 0.207 & 1.57 & 0.00955 & 0.0509 & 0.874 & 3.17 & 0.272 & 0.288 \\
ReLU & active & 1 & 0.786 & 0.51 & 1.57 & 0.00955 & 0.0509 & 1.73 & 3.17 & 0.501 & 0.301 \\
ReLU & active & 4 & 0.765 & 0.815 & 1.57 & 0.00955 & 0.043 & 2.43 & 3.17 & 0.558 & 0.338 \\
ReLU & inactive & 0.25 & 1.69 & 0.798 & 4.06 & 0.115 & 0.218 & 3.14 & 7.63 & 0.84 & 0.725 \\
ReLU & inactive & 1 & 0.862 & 0.803 & 2.26 & 0.0792 & 0.127 & 2.41 & 4.12 & 0.746 & 0.343 \\
ReLU & inactive & 4 & 0.603 & 1.07 & 1.82 & 0.0674 & 0.0931 & 2.38 & 3.51 & 0.653 & 0.307 \\
\midrule
JumpReLU & active & 0.25 & 1.18 & 1.6 & 1.54 & $<\!10^{-6}$ & $<\!10^{-6}$ & 2.54 & 3.22 & 0.227 & 0.892 \\
JumpReLU & active & 1 & 0.928 & 1.2 & 1.54 & $<\!10^{-6}$ & $<\!10^{-6}$ & 2.63 & 3.22 & 0.487 & 0.701 \\
JumpReLU & active & 4 & 0.855 & 1.02 & 1.54 & $<\!10^{-6}$ & $<\!10^{-6}$ & 2.59 & 3.22 & 0.546 & 0.499 \\
JumpReLU & inactive$^\dagger$ & 0.25 & 1.53 & 1.51 & 1.96 & 1 & 0.0306 & 2.66 & 3.54 & 1 & 0.903 \\
JumpReLU & inactive & 1 & 0.935 & 1.25 & 1.72 & 0.00982 & 0.0803 & 2.57 & 3.06 & 0.575 & 0.68 \\
JumpReLU & inactive & 4 & 0.611 & 1.25 & 1.65 & 0.0101 & 0.0512 & 2.42 & 3 & 0.597 & 0.473 \\
\midrule
TopK & active & 0.25 & 1.39 & 0.993 & 1.59 & $<\!10^{-6}$ & 1.36 & 1.57 & 2.11 & $<\!10^{-6}$ & 0.206 \\
TopK & active & 1 & 0.942 & 0.837 & 1.59 & $<\!10^{-6}$ & 0.865 & 1.35 & 2.11 & 0.511 & 2.38 \\
TopK & active & 4 & 0.707 & 0.86 & 1.59 & $<\!10^{-6}$ & 0.364 & 1.32 & 2.11 & 0.394 & 0.794 \\
TopK & inactive & 0.25 & 1.36 & 1.08 & 1.66 & 1 & 1.08 & 1.4 & 2.07 & 1 & 0.603 \\
TopK & inactive & 1 & 0.98 & 0.878 & 1.59 & $<\!10^{-6}$ & 0.944 & 1.28 & 2.02 & 0.465 & 1.67 \\
TopK & inactive & 4 & 0.55 & 0.96 & 1.64 & $<\!10^{-6}$ & 0.4 & 1.37 & 2.07 & 0.432 & 0.745 \\
\bottomrule
\end{tabular}

\end{table}

\begin{table}[h]
  \caption{Exploratory sparser ReLU settings (added after development seeds).}
  \label{tab:sparse}
  \centering
  \small
  \begin{tabular}{llc|ccc|cccc|cc}
\toprule
 & & & \multicolumn{3}{c|}{Rescaled decoder} & \multicolumn{4}{c|}{LN correction} & \multicolumn{2}{c}{LN + repair} \\
Act. & Target & $\alpha$ & orig.\ act. & new & norm & tgt err & orig.\ act. & new & norm & tgt err & total \\
\midrule
$b{=}{-}0.5$ & active & 0.25 & 0.562 & 0.23 & 1.5 & $<\!10^{-6}$ & $<\!10^{-6}$ & 0.244 & 1.86 & $<\!10^{-6}$ & $<\!10^{-6}$ \\
$b{=}{-}0.5$ & active & 1 & 0.561 & 0.507 & 1.5 & $<\!10^{-6}$ & $<\!10^{-6}$ & 0.601 & 1.86 & 0.403 & 0.216 \\
$b{=}{-}0.5$ & active & 4 & 0.558 & 0.892 & 1.5 & $<\!10^{-6}$ & $<\!10^{-6}$ & 1.16 & 1.86 & 0.515 & 0.269 \\
$b{=}{-}0.5$ & inactive & 0.25 & 1.09 & 0.689 & 4.71 & $<\!10^{-6}$ & $<\!10^{-6}$ & 1.04 & 5.42 & 0.495 & 0.884 \\
$b{=}{-}0.5$ & inactive & 1 & 0.498 & 0.77 & 2.35 & $<\!10^{-6}$ & $<\!10^{-6}$ & 0.92 & 2.79 & 0.67 & 0.24 \\
$b{=}{-}0.5$ & inactive & 4 & 0.343 & 1.09 & 1.78 & $<\!10^{-6}$ & $<\!10^{-6}$ & 1.29 & 2.13 & 0.62 & 0.234 \\
\midrule
$b{=}{-}0.75$ & active & 0.25 & 0.137 & 0.0905 & 1.34 & $<\!10^{-6}$ & $<\!10^{-6}$ & 0.0693 & 1.38 & $<\!10^{-6}$ & $<\!10^{-6}$ \\
$b{=}{-}0.75$ & active & 1 & 0.137 & 0.391 & 1.34 & $<\!10^{-6}$ & $<\!10^{-6}$ & 0.351 & 1.38 & 1.9e-04 & 0.0976 \\
$b{=}{-}0.75$ & active & 4 & 0.137 & 0.789 & 1.34 & $<\!10^{-6}$ & $<\!10^{-6}$ & 0.828 & 1.38 & 0.484 & 0.18 \\
$b{=}{-}0.75$ & inactive & 0.25 & 0.502 & 0.616 & 6.67 & $<\!10^{-6}$ & $<\!10^{-6}$ & 0.68 & 6.8 & $<\!10^{-6}$ & 0.408 \\
$b{=}{-}0.75$ & inactive & 1 & 0.212 & 0.628 & 2.87 & $<\!10^{-6}$ & $<\!10^{-6}$ & 0.689 & 2.94 & 0.545 & 0.493 \\
$b{=}{-}0.75$ & inactive & 4 & 0.124 & 1.04 & 1.94 & $<\!10^{-6}$ & $<\!10^{-6}$ & 1.07 & 2.02 & 0.635 & 0.128 \\
\bottomrule
\end{tabular}

\end{table}

\begin{table*}[h]
  \caption{Pre-registered dense ReLU setting (20 instances; medians). Drift is relative to
  $\alpha$. ``orig.\ act.'' is $D_{\mathrm{orig}}$ and ``new'' is $D_{\mathrm{new}}$
  (\cref{def:fidelity}). The last two columns are target error and total drift after
  repair. Other activations are in \cref{tab:activeset-all}.}
  \label{tab:relu}
  \centering
  \small
  \setlength{\tabcolsep}{3.5pt}
  \begin{tabular}{llc|ccc|cccc|cc}
\toprule
 & & & \multicolumn{3}{c|}{Rescaled decoder} & \multicolumn{4}{c|}{LN correction} & \multicolumn{2}{c}{LN + repair} \\
Act. & Target & $\alpha$ & orig.\ act. & new & norm & tgt err & orig.\ act. & new & norm & tgt err & total \\
\midrule
ReLU & active & 0.25 & 0.827 & 0.207 & 1.57 & 0.00955 & 0.0509 & 0.874 & 3.17 & 0.272 & 0.288 \\
ReLU & active & 1 & 0.786 & 0.51 & 1.57 & 0.00955 & 0.0509 & 1.73 & 3.17 & 0.501 & 0.301 \\
ReLU & active & 4 & 0.765 & 0.815 & 1.57 & 0.00955 & 0.043 & 2.43 & 3.17 & 0.558 & 0.338 \\
ReLU & inactive & 0.25 & 1.69 & 0.798 & 4.06 & 0.115 & 0.218 & 3.14 & 7.63 & 0.84 & 0.725 \\
ReLU & inactive & 1 & 0.862 & 0.803 & 2.26 & 0.0792 & 0.127 & 2.41 & 4.12 & 0.746 & 0.343 \\
ReLU & inactive & 4 & 0.603 & 1.07 & 1.82 & 0.0674 & 0.0931 & 2.38 & 3.51 & 0.653 & 0.307 \\
\bottomrule
\end{tabular}

\end{table*}

\begin{table}[h]
  \caption{Synthetic external proxy (20 instances; medians). ``Err'' is
  $\|\Delta z - \alpha e_j\|/\alpha$ measured with the SAE encoder (internal) or with $A^{+}$
  (external). The last column counts instances where the \LN{} correction beats the decoder
  externally.}
  \label{tab:proxy}
  \centering
  \scriptsize
  \setlength{\tabcolsep}{3pt}
  \begin{tabular}{llcccc}
\toprule
Case & Method & Internal err & External err & Norm & LN better ext. \\
\midrule
encoder $=A^{+}$ & decoder & 0.221 & 0.221 & 1 &  \\
 & decoder\_rescaled & 0.222 & 0.222 & 1.04 &  \\
 & encoder\_grad & 0.283 & 0.283 & 0.88 &  \\
 & jacobian\_ln & $<\!10^{-6}$ & $<\!10^{-6}$ & 1 & 20/20 \\
\midrule
decoder $=A$ & decoder & 0.242 & $<\!10^{-6}$ & 1 &  \\
 & decoder\_rescaled & 0.226 & 0.0708 & 0.997 &  \\
 & encoder\_grad & 0.362 & 0.33 & 0.846 &  \\
 & jacobian\_ln & $<\!10^{-6}$ & 0.239 & 0.964 & 0/20 \\
\midrule
both noisy & decoder & 0.313 & 0.223 & 1 &  \\
 & decoder\_rescaled & 0.306 & 0.241 & 1.06 &  \\
 & encoder\_grad & 0.367 & 0.334 & 0.845 &  \\
 & jacobian\_ln & $<\!10^{-6}$ & 0.223 & 0.986 & 11/20 \\
\bottomrule
\end{tabular}

\end{table}

\section{Model--SAE Contract}
\label{app:contract}
\begin{table}[h]
  \caption{Contract fields, sources and status after execution. ``Verified'' means read from
  the pinned source file or code; ``executed'' means checked by running both paths
  (\cref{tab:contract-exec}).}
  \label{tab:contract}
  \centering
  \small
  \begin{tabular}{p{0.24\linewidth}p{0.42\linewidth}p{0.24\linewidth}}
    \toprule
    Field & Value & Status \\
    \midrule
    Model & GPT-2 small, pinned HF revision, MIT (card) & verified \\
    SAE & residual SAE, \texttt{blocks.8.hook\_resid\_pre}, pinned revision, MIT (card) & verified; weights bitwise equal to the SAELens load \\
    Width, dtype, context & $d=768$, $m=24576$, float32 weights, 128 & verified (config, header) \\
    Activation, input bias, normalization & ReLU; $x - b_{\mathrm{dec}}$; none & library defaults, executed (K0, K3) \\
    Coordinates & HF hidden state minus per-token mean & executed (K2, K4) \\
    Tensors & $W_{\mathrm{enc}}\,[768{\times}24576]$, $W_{\mathrm{dec}}\,[24576{\times}768]$, biases, F32 & verified (header) \\
    Tokenizer, BOS, positions & GPT-2 BPE, BOS prepended, absolute positions & executed (K1) \\
    Processing options & fold\_ln, center\_writing\_weights, center\_unembed, fold\_value\_biases on (explicit) & executed; no-processing path agrees (K10) \\
    Model-visible edits & classes $\delta+\operatorname{span}(\mathbf{1})$, represented in $\operatorname{range}(P)$, $P = I - \mathbf{1}\mathbf{1}^{\top}/d$ & proved for GPT-2 (\cref{prop:admissible}); executed (K9) \\
    SAE readout of an edit & canonical $E(P(h+\delta)) - E(Ph)$ & readout closure (\cref{sec:coord}) \\
    Data & WikiText-103 raw (CC BY-SA 3.0 / GFDL on card) & verified (card) \\
    Libraries & TransformerLens 2.18.0, SAELens 6.51.3, torch 2.14.0 (CPU) & pinned \\
    \bottomrule
  \end{tabular}
\end{table}

\section{Real-Model Pilot Details}
\label{app:pilot}
\paragraph{Frozen choices.} All of the following were fixed from calibration documents only
and frozen, with a hash recorded, before any test document was read:
\begin{itemize}
  \item the calibration feature pool (firing density in $[10^{-4},10^{-2}]$);
  \item the 64 sampled features and the pilot subsample of \PilotFeatures;
  \item per-feature doses ($q_{50}$, $q_{90}$, $q_{99}$, $2q_{99}$ of positive calibration
  activations; $a_{\max}$ recorded but unused);
  \item norm budgets ($\{0.025, 0.05, 0.1, 0.2\}$ times the median centred residual norm).
\end{itemize}

\paragraph{Changes before and during the test run.} A timing smoke on calibration documents
projected an upper bound that exceeded the 30-minute budget, because its formula counted
model loading and scanning once per group. Before the test run, we kept the pre-registered
pilot size ($N = 8$ documents per feature and target type) after estimating the cost per
quality group from the smoke's own rows. Two rules were committed to the summary code while
the test run was executing and before any test row existed on disk. They are \emph{in-run
amendments}, not pre-registered:
\begin{itemize}
  \item gate readouts need at least 10 documents;
  \item a derived target-gain metric is reported.
\end{itemize}

\paragraph{Gate readouts (descriptive).} The earlier gate conditions G1--G3 were read at
$q_{99}$ with document-cluster intervals, and all three fell on the continue side:
\begin{itemize}
  \item G1: rescaled originally-active drift is above $0.1\alpha$. As run, it pools matched
  (active) and unmatched (inactive) rows; the split is in \cref{sec:denominators};
  \item G2: newly activated drift is below originally-active drift;
  \item G3: the correction lowers total drift.
\end{itemize}
The pilot configuration declares these readouts descriptive, not decisions.

\begin{table*}[h]
  \caption{Pilot, target-matched at $q_{99}$: medians with 95\% document-cluster bootstrap
  intervals. $^{k}$: $k$ infeasible rows (kept, not dropped).}
  \label{tab:pilot-ci}
  \centering
  \tiny
  \setlength{\tabcolsep}{2pt}
  \begin{tabular}{llcccc}
\toprule
Target & Method & Target err & $D_{\mathrm{orig}}/\alpha$ & $D_{\mathrm{new}}/\alpha$ & $D_{\mathrm{all}}/\alpha$ \\
\midrule
active & decoder & $<10^{-7}$ & 0.583\,\allowbreak[0.452, 0.6] & 0.311\,\allowbreak[0.236, 0.342] & 0.656\,\allowbreak[0.55, 0.704] \\
 & rescaled (weight-only) & $<10^{-7}$ & 0.583\,\allowbreak[0.452, 0.6] & 0.311\,\allowbreak[0.236, 0.342] & 0.656\,\allowbreak[0.55, 0.704] \\
 & encoder row & $<10^{-7}$ & 0.681\,\allowbreak[0.608, 0.725] & 0.366\,\allowbreak[0.314, 0.45] & 0.778\,\allowbreak[0.691, 0.935] \\
 & LN ($J_EP$) & $<10^{-7}$ & \ensuremath{4.7\times10^{-15}}\,\allowbreak[\ensuremath{3.6\times10^{-15}}, \ensuremath{5.7\times10^{-15}}] & 0.114\,\allowbreak[0.0976, 0.135] & 0.114\,\allowbreak[0.0976, 0.135] \\
 & random$^{33}$ & $<10^{-7}$ & 7.33\,\allowbreak[4.32, 10.9] & 108\,\allowbreak[46.8, 237] & 108\,\allowbreak[46.9, 238] \\
\midrule
inactive & decoder & $<10^{-7}$ & 0.583\,\allowbreak[0.485, 0.685] & 0.361\,\allowbreak[0.303, 0.416] & 0.716\,\allowbreak[0.59, 0.819] \\
 & rescaled (weight-only) & 0.727\,\allowbreak[0.629, 0.856] & 0.342\,\allowbreak[0.289, 0.371] & 0.112\,\allowbreak[0.0993, 0.125] & 0.37\,\allowbreak[0.322, 0.394] \\
 & encoder row & $<10^{-7}$ & 0.943\,\allowbreak[0.707, 1.2] & 0.943\,\allowbreak[0.784, 1.07] & 1.36\,\allowbreak[1.13, 1.53] \\
 & LN ($J_EP$) & $<10^{-7}$ & \ensuremath{5.0\times10^{-15}}\,\allowbreak[\ensuremath{4.0\times10^{-15}}, \ensuremath{6.5\times10^{-15}}] & 0.295\,\allowbreak[0.27, 0.342] & 0.295\,\allowbreak[0.27, 0.342] \\
 & random$^{35}$ & $<10^{-7}$ & 17.4\,\allowbreak[7.69, 32] & 347\,\allowbreak[222, 604] & 348\,\allowbreak[222, 604] \\
\bottomrule
\end{tabular}

\end{table*}

\begin{table*}[h]
  \caption{Pilot: per-document paired differences, \LN{} minus decoder (target-matched), with
  95\% document-cluster bootstrap intervals.}
  \label{tab:pilot-paired}
  \centering
  \scriptsize
  \begin{tabular}{llccccc}
\toprule
Target & Dose & $n_{\mathrm{doc}}$ & $\Delta D_{\mathrm{all}}$ & $\Delta D_{\mathrm{new}}$ & $\Delta D_{\mathrm{orig}}$ & $\Delta$KL \\
\midrule
active & q50 & 42 & -0.542 [-0.586, -0.455] & -0.0337 [-0.0499, -0.0238] & -0.588 [-0.607, -0.459] & -1.1e-05 [-3.7e-05, -6.4e-06] \\
 & q90 & 42 & -0.52 [-0.577, -0.438] & -0.0927 [-0.115, -0.0686] & -0.587 [-0.603, -0.456] & -2.3e-04 [-7.2e-04, -4.4e-05] \\
 & q99 & 42 & -0.503 [-0.584, -0.421] & -0.17 [-0.202, -0.142] & -0.583 [-0.6, -0.452] & -0.00377 [-0.0117, -2.5e-04] \\
 & 2xq99 & 42 & -0.518 [-0.596, -0.421] & -0.282 [-0.304, -0.216] & -0.566 [-0.599, -0.451] & -0.0121 [-0.0502, -8.8e-04] \\
\midrule
inactive & q50 & 40 & -2.58 [-2.97, -2.24] & -0.119 [-0.259, 0.0136] & -3.18 [-3.47, -2.89] & -4.0e-04 [-6.3e-04, -1.0e-04] \\
 & q90 & 40 & -0.741 [-0.925, -0.597] & -0.0729 [-0.0909, 0.0173] & -0.973 [-1.15, -0.876] & -7.3e-04 [-0.00149, -3.2e-04] \\
 & q99 & 40 & -0.43 [-0.517, -0.35] & -0.0557 [-0.098, -0.00262] & -0.583 [-0.685, -0.485] & -0.00234 [-0.00651, -8.7e-04] \\
 & 2xq99 & 40 & -0.294 [-0.348, -0.25] & -0.0467 [-0.133, -0.0127] & -0.423 [-0.507, -0.372] & -0.00589 [-0.0212, -0.00205] \\
\bottomrule
\end{tabular}

\end{table*}

\begin{table}[h]
  \caption{Pilot, equal norm ($\rho = 0.1\times$ median residual norm). ``Target gain'' is the
  target-feature change per unit edit norm, a metric added as an in-run amendment
  (\cref{sec:denominators}). Mean-only rows are omitted (raw off-slice readout).}
  \label{tab:pilot-eq}
  \centering
  \scriptsize
  \setlength{\tabcolsep}{2pt}
  \begin{tabular}{llccc}
\toprule
Target & Method & target gain & $D_{\mathrm{all}}/\rho$ & KL (nats) \\
\midrule
active & decoder & 0.895 [0.847, 0.946] & 0.517 [0.473, 0.53] & 7.3e-04 [4.2e-04, 0.00159] \\
 & encoder row & 1.5 [1.47, 1.72] & 1.16 [0.967, 1.26] & 8.0e-04 [5.1e-04, 0.00119] \\
 & LN ($J_EP$) & 1.35 [1.29, 1.55] & 0.142 [0.126, 0.152] & 6.5e-04 [3.6e-04, 8.4e-04] \\
 & random & 0.0383 [0.029, 0.0483] & 0.341 [0.319, 0.361] & 4.6e-04 [3.8e-04, 4.8e-04] \\
\midrule
inactive & decoder & 0 [0, 0.0494] & 0.316 [0.303, 0.335] & 5.2e-04 [2.4e-04, 6.3e-04] \\
 & encoder row & 0.433 [0.303, 0.476] & 0.859 [0.726, 0.984] & 5.2e-04 [3.2e-04, 6.8e-04] \\
 & LN ($J_EP$) & 0.304 [0.174, 0.386] & 0.125 [0.101, 0.141] & 4.4e-04 [2.9e-04, 6.3e-04] \\
 & random & 0 [0, 0] & 0.348 [0.32, 0.38] & 2.7e-04 [2.3e-04, 5.9e-04] \\
\bottomrule
\end{tabular}

\end{table}

\section{P3-REAL-01R Protocol}
\label{app:real01r}
\paragraph{Changes from the earlier design.}
\begin{itemize}
  \item The unit is the document, not the feature.
  \item $a_{\max}$ is replaced by calibration-only quantiles.
  \item Features and doses are frozen before any test window is encoded.
  \item Equal-norm and target-matched comparisons are reported separately.
  \item Infeasible rows are kept.
\end{itemize}

\paragraph{Gates.}
The original gate text is preserved:
\begin{itemize}
  \item G1: ``the rescaled decoder's protected leakage is $<0.1\alpha$'';
  \item G2: ``crossing leakage $\ge$ protected leakage for the rescaled decoder'';
  \item G3: ``the correction's total leakage $\ge$ the rescaled decoder's in $\ge 50\%$ of
  features''.
\end{itemize}
The condition used to be the median over test features at $\alpha\ge a_{\max}$. It is now
read at $q_{99}$, per document, with 95\% document-cluster bootstrap intervals, and each
readout is one of stop side, continue side, or inconclusive. A large crossing share ends
the method claim but not the measurement study.

\paragraph{Stages.}
Six command-line stages record their status (ran, blocked, not run) with time and peak
memory: environment check, contract, calibration, timing smoke on calibration documents
only, run, and summary. All six stages ran for the pilot (\cref{app:pilot}); the full run over all 64
calibration features is \notrun.

\section{P3-REAL-02-LX Protocol}
\label{app:real02}
The configuration was committed before the calibration stage and before any test document
was read.
\begin{itemize}
  \item \emph{Metric.} Wedding-lexicon occurrence in the 32 generated tokens (prompt
  excluded), using the ActAdd keyword list \citep{turner2024steering}: \emph{wedding,
  weddings, wed, marry, married, marriage, bride, groom, honeymoon}. We report the share of
  documents with at least one hit and the hits per generated token.
  \item \emph{Calibration only (WikiText-103 validation).} The feature is the one with the
  largest activation difference between windows with and without keywords, and the
  DiffMean direction is the difference of mean centred hook states. The lexicon is thus used
  in direction selection. Each method's norm budget is the largest of
  $\{0.05, 0.1, 0.2, 0.4, 0.8\}$ times the median centred norm whose mean teacher-forced
  KL$(P_{\mathrm{base}}\,\|\,P_{\mathrm{edit}})$ on 32 calibration prompts is at most
  $\kappa = 0.1$ nats; a method with no such budget is \notrun, not substituted.
  \item \emph{Methods.} Projected decoder $PD_j$, projected encoder row $PE_j^{\top}$,
  per-position \LN{} correction ($J_EP$), projected DiffMean, and no edit. Random directions,
  SAE-TS and FGAA are \notrun. The mean-only control runs at the largest budget, outside
  every ranking.
  \item \emph{Test.} Test articles sorted by identifier and shuffled with a fixed seed; the
  first 32 (or fewer, by a cost rule fixed in advance) with a keyword-free window give one
  prompt each (BOS plus 31 tokens). Greedy decoding (temperature 0, no sampling), 32 new
  tokens, edit at every position, no KV cache.
  \item \emph{Quality.} KL on the common prompt (teacher forced, averaged over positions) and
  the continuation's NLL under the unedited model. Internal: canonical target change and
  non-target drift at prompt positions, never used as the label.
  \item \emph{Analysis.} The unit is the test document; means over documents with
  document-bootstrap 95\% intervals. The method claim is supported only if \LN{} beats
  decoder, encoder row and DiffMean with intervals excluding zero.
  \item \emph{Amendment v1.1.} The first calibration (windows capped at 8 per document) found
  only 5 keyword windows, below the pre-set minimum of 20, and returned \textsc{blocked}. We
  kept that record and committed one amendment before reading any test document: scan every
  calibration window, everything else unchanged; a second \textsc{blocked} would have ended
  the experiment.
\end{itemize}

\section{Reproducibility, Compute and Claim Map}
\label{app:repro}
\begin{itemize}
  \item \emph{Toy suite.} It runs in 120\,s or less in the standard library from a single
  configuration. Tables, figure data and every number in the text are regenerated by one
  script.
  \item \emph{Real-model runs.} They used pinned revisions and a pinned environment
  (package list in the results), in float64 on two CPU threads. Wall-clock times:
  \begin{itemize}
    \item contract execution: 161\,s in float32 and 391\,s in float64;
    \item adapter contract stage: 71\,s; calibration: 330\,s; timing smoke: 95\,s;
    pilot run: 1043\,s; summary: 34\,s;
    \item readout closure: \ClWall\,s wall (\ClCPU{} CPU-seconds summed over the two threads);
    \item matched-denominator re-aggregation: about 5\,s (standard library, one thread);
    \item REAL-02-LX: adapter checks 9\,s; calibration \LxCalSecVOne\,s (v1, blocked) and
    \LxCalSecVOneOne\,s (v1.1, blocked); run and summary not executed.
  \end{itemize}
  Peak memory was below 9\,GB. No GPU was used, and nothing was installed or downloaded
  after the third session.
  \item \emph{Claim map.} A claim-to-evidence table (\texttt{claims.csv}) links each
  statement to its files, experiment identifiers, assumptions and status.
  \item \emph{Release.} The code and raw files will be released in anonymized form.
\end{itemize}

\end{document}

